
%%FILE numina/blueprints/16-made-easy/16-made-easy.tex
% Blueprint: 16 made easy
% Created by Numina

\begin{document}

\section{Setup}

We follow Wild's paper ``The Groups of Order Sixteen Made Easy'' (Marcel Wild,
\emph{Amer.\ Math.\ Monthly} \textbf{112} (2005), 20--31). Throughout, $G$ is a
finite group of order $16$ and is not isomorphic to $(C_2)^4$. The goal is to
show that $G$ realises one of the $13$ cyclic-extension types
$\mathtt{ext\_16\_1}, \dots, \mathtt{ext\_16\_13}$ already defined in
\texttt{M2rGroup7/SixteenClassification/Blueprints.lean}.

We use the following data structures from
\texttt{M2rGroup7/SixteenClassification/Extensions.lean}:

\begin{itemize}
\item An \emph{extension type} $E = (N, n, \tau, v)$ packages a group $N$, an
exponent $n \in \mathbb{N}$, an automorphism $\tau \in \mathrm{Aut}(N)$ fixing
$v \in N$, with $\tau^n = \mathrm{conj}_v$.
\item A \emph{realisation} of $E$ by a group $G$ is a tuple
$(a, \iota)$ where $a \in G$, $\iota : N \hookrightarrow G$ is a monoid
homomorphism, $a \iota(x) a^{-1} = \iota(\tau(x))$ for all $x \in N$,
$a^n = \iota(v)$, and the map $(x, i) \mapsto \iota(x) \cdot a^i$ is a bijection
$N \times \{0, 1, \dots, n-1\} \xrightarrow{\sim} G$.
\end{itemize}

The thirteen extension types $\mathtt{ext\_16\_i}$ are pre-defined; each one is
realised by an explicit order-$16$ group in
\texttt{M2rGroup7/SixteenClassification/Blueprints.lean}.

\section{Transfer along isomorphisms}

\begin{definition}
  \label{def:realise-transfer}
  \lean{OrderSixteen.RealiseExtType.transfer}
  \leanfile{M2rGroup7/SixteenClassification/Lemma3.lean}
  \leanok
  Let $G$ and $G'$ be groups, $E$ an extension type, $e : G \simeq G'$ a group
  isomorphism, and $R : \mathrm{RealiseExtType}~G'~E$. Then the
  \emph{transferred realisation} $\mathrm{transfer}(e, R)$ is the canonical
  realisation of $E$ by $G$ obtained by pulling $R$ back along $e$.
  Concretely, take $a := e^{-1}(R.a)$, $\iota := e^{-1} \circ R.\iota$,
  and the bijection $(x, i) \mapsto \iota(x) \cdot a^i$ as $R.\mathrm{equiv}$
  composed with $e^{-1}$.
\end{definition}

\begin{lemma}
  \label{lem:realise-of-mulEquiv}
  \lean{OrderSixteen.realise_of_mulEquiv}
  \leanfile{M2rGroup7/SixteenClassification/Lemma3.lean}
  \uses{def:realise-transfer}
  \leanok
  If $G \simeq G'$ and $G'$ admits a realisation of an extension type $E$, then
  so does $G$.
\end{lemma}

\begin{proof}
  Immediate from Definition~\ref{def:realise-transfer}.
\end{proof}

\section{Inducing elements}

\begin{lemma}
  \label{lem:inducing-element-exists}
  \lean{OrderSixteen.exists_inducing_element}
  \leanfile{M2rGroup7/SixteenClassification/Lemma3.lean}
  Let $G$ be a group and $H \triangleleft G$ a normal subgroup of index $2$.
  Then there exists $a \in G \setminus H$, and any such $a$ satisfies
  $a^2 \in H$ and $\mathrm{conj}_a(H) = H$. The map $h \mapsto a h a^{-1}$ is
  therefore an automorphism of $H$.
\end{lemma}

\begin{proof}
  Since $[G : H] = 2 > 1$, $H \neq G$, so some $a \notin H$ exists. The coset
  $Ha$ has order dividing $[G : H] = 2$ and is non-identity, hence order $2$.
  Therefore $(Ha)^2 = H$, i.e.\ $a^2 \in H$. Normality of $H$ gives
  $aHa^{-1} = H$ and the conjugation map $h \mapsto aha^{-1}$ restricts to an
  automorphism of $H$.
\end{proof}

\begin{lemma}
  \label{lem:realise-from-normal-index-2}
  \lean{OrderSixteen.realise_from_normal_index_two}
  \leanfile{M2rGroup7/SixteenClassification/Lemma3.lean}
  \uses{lem:inducing-element-exists}
  Let $G$ be a finite group, $H \triangleleft G$ a normal subgroup of index $2$,
  and $a \in G \setminus H$. Set $v := a^2 \in H$ and let
  $\tau \in \mathrm{Aut}(H)$ be conjugation by $a$. Then the extension type
  $E := (H, 2, \tau, v)$ is well-defined ($\tau(v) = v$ and $\tau^2 =
  \mathrm{conj}_v$) and $G$ realises $E$ with inducing element $a$ and
  embedding $\iota : H \hookrightarrow G$.
\end{lemma}

\begin{proof}
  Set $v := a^2$ and $\tau(h) := aha^{-1}$. Then $\tau(v) = a v a^{-1} = a^3
  a^{-1} = a^2 = v$ and $\tau^2(h) = a^2 h a^{-2} = v h v^{-1} = \mathrm{conj}_v
  (h)$. The decomposition $G = H \sqcup Ha$ as cosets of $H$ gives the
  bijection $H \times \{0, 1\} \xrightarrow{\sim} G$,
  $(h, i) \mapsto h \cdot a^i$, completing the data of a realisation.
\end{proof}

\section{Case analysis for normal \texorpdfstring{$C_8$}{C8}}

\begin{lemma}
  \label{lem:realise-with-normal-C8}
  \lean{OrderSixteen.realise_with_normal_C8}
  \leanfile{M2rGroup7/SixteenClassification/Lemma3.lean}
  \uses{lem:realise-from-normal-index-2, lem:realise-of-mulEquiv}
  Let $G$ be a group of order $16$ with a normal subgroup $H \le G$ such that
  $H \simeq C_8$. Then $G$ realises one of the extension types
  $\mathtt{ext\_16\_1}$, $\mathtt{ext\_16\_5}$, $\mathtt{ext\_16\_6}$,
  $\mathtt{ext\_16\_7}$, $\mathtt{ext\_16\_8}$, or $\mathtt{ext\_16\_9}$.
\end{lemma}

\begin{proof}
  Write $H = \langle x \rangle$ with $|x| = 8$. The hypotheses $|G| = 16$ and
  $|H| = 8$ give $[G : H] = 2$, so $H$ is normal of index $2$ and Lagrange
  gives a partition $G = H \sqcup Ha$ for any chosen
  $a \in G \setminus H$.

  By Lemma~\ref{lem:inducing-element-exists} the set $G \setminus H$ is
  non-empty, so the set
  $\{\,|b| \mid b \in G \setminus H\,\}$ is a non-empty subset of the positive
  divisors of $|G| = 16$, hence has a minimum. Choose an inducing element
  $a \in G \setminus H$ of \emph{minimal order}. Because $a \notin H$ but
  $a^2 \in H$ (since the coset $Ha$ has order $2$ in $G/H$), the order of $a$
  is even, and as $|a|$ divides $|G| = 16$ we have
  $|a| \in \{2, 4, 8, 16\}$.

  Set $v := a^2 \in H$ and let $\tau := \mathrm{conj}_a \in \mathrm{Aut}(H)$ be
  conjugation by $a$ restricted to $H$. By
  Lemma~\ref{lem:realise-from-normal-index-2}, $G$ realises the extension type
  $(H, 2, \tau, v)$.

  Recall (Wild's Fact~3 for $\mathrm{Aut}(C_8)$) that
  $\mathrm{Aut}(H) \cong (\mathbb{Z}/8\mathbb{Z})^\times =
  \{1, 3, 5, 7\}$ as a multiplicative group, with $\tau$ acting as $x \mapsto
  x^k$ for one of the four odd residues $k$. Every such $k$ satisfies
  $k^2 \equiv 1 \pmod 8$, so every $\tau \in \mathrm{Aut}(H)$ is an involution.

  We split on $|a|$.

  \smallskip
  \noindent
  \emph{Case 1: $|a| = 2$.} Then $v = a^2 = e$, so $\mathrm{conj}_v$ is the
  identity automorphism of $H$, and the validity condition $\tau^2 =
  \mathrm{conj}_v$ becomes $\tau^2 = \mathrm{id}$. This is automatic since every
  $\tau \in \mathrm{Aut}(H)$ is an involution. The four choices
  $\tau \in \{1, \cdot^3, \cdot^5, \cdot^7\}$ give four extension types
  $(H, 2, \tau, e)$, matching exactly the four predefined types
  \[
    \mathtt{ext\_16\_5} = (C_8, 2, \mathrm{id}, e),\quad
    \mathtt{ext\_16\_8} = (C_8, 2, \cdot^3, e),\quad
    \mathtt{ext\_16\_6} = (C_8, 2, \cdot^5, e),\quad
    \mathtt{ext\_16\_7} = (C_8, 2, \cdot^7, e),
  \]
  realised respectively by $C_8 \times C_2$, $C_8 \rtimes_{\mathrm{pow}_3} C_2$,
  $C_8 \rtimes_{\mathrm{pow}_5} C_2$, and the dihedral group $D_8$. In each
  sub-case the realisation of $G$ is obtained directly from the inducing data
  $(a, \iota)$ produced by
  Lemma~\ref{lem:realise-from-normal-index-2}.

  \smallskip
  \noindent
  \emph{Case 2: $|a| = 4$.} Then $v = a^2$ has order $2$, and the only
  involution in $H = \langle x \rangle \simeq C_8$ is $x^4$, so $v = x^4$.
  Minimality means no element of $G \setminus H$ has order $2$, otherwise we
  would be in Case~1. The validity constraint $\tau(v) = v$ is automatic:
  $\tau(x^4) = x^{4k} = x^4$ for all odd $k$ (since $4k \equiv 4 \pmod 8$).
  The validity constraint $\tau^2 = \mathrm{conj}_{x^4} = \mathrm{id}$ (as
  $H$ is abelian) is also automatic. So all four
  $\tau \in \{1, \cdot^3, \cdot^5, \cdot^7\}$ pass the validity checks, giving
  candidate extension types $(C_8, 2, \tau, x^4)$. We show that the minimality
  assumption rules out three of them:
  \begin{itemize}
    \item $\tau = \mathrm{id}$: then $a x = x a$, and $b := x^{-2} a$
    satisfies $b^2 = x^{-4} a^2 = x^{-4} x^4 = e$, so $b$ is an involution. As
    $a \notin H$ and $x^{-2} \in H$, $b \notin H$, contradicting minimality.
    \item $\tau = \cdot^3$: then $a x a^{-1} = x^3$, and $(xa)^2 = x \cdot
    a x a^{-1} \cdot a^2 = x \cdot x^3 \cdot x^4 = x^8 = e$. Thus $xa$ is an
    involution in $G \setminus H$, contradicting minimality.
    \item $\tau = \cdot^5$: then $a x^2 a^{-1} = (a x a^{-1})^2 = x^{10}$, and
    $(x^2 a)^2 = x^2 \cdot a x^2 a^{-1} \cdot a^2 = x^2 \cdot x^{10} \cdot x^4
    = x^{16} = e$. Thus $x^2 a$ is an involution in $G \setminus H$,
    contradicting minimality.
    \item $\tau = \cdot^7$ (i.e.\ inversion): for every $k$,
    $(x^k a)^2 = x^k \cdot a x^k a^{-1} \cdot a^2 = x^k \cdot x^{-k} \cdot x^4
    = x^4 \neq e$, and $(x^k a)^4 = (x^4)^2 = x^8 = e$, so every element of
    $G \setminus H$ has order $4$. This satisfies the minimality assumption.
  \end{itemize}
  Hence only $\tau = \cdot^7$ survives in Case~2, yielding the extension type
  $\mathtt{ext\_16\_9} = (C_8, 2, \cdot^7, x^4)$. The corresponding realising
  group is the generalised quaternion group $Q_{16}$.

  \smallskip
  \noindent
  \emph{Case 3: $|a| = 8$.} Then $v = a^2$ has order $4$, so
  $v \in \{x^2, x^6\}$. The two values are related by $x^6 = (x^2)^3$, and
  the automorphism $\cdot^3$ sends $x^2$ to $x^6$. The action set
  $\{\mathrm{id}, \cdot^3, \cdot^5, \cdot^7\}$ is a union of
  $\mathrm{Aut}(H)$-conjugacy classes (every element is in its own class
  because $\mathrm{Aut}(H)$ is abelian, but this is enough). Hence by
  Lemma~\ref{lem:realise-of-mulEquiv} together with
  \texttt{extension\_families\_same\_isoClasses} we may assume $v = x^2$.

  The validity constraint $\tau(v) = v$ now reads $x^{2k} = x^2$, i.e.\
  $2(k - 1) \equiv 0 \pmod 8$, i.e.\ $k \equiv 1 \pmod 4$, so
  $k \in \{1, 5\}$ and $\tau \in \{\mathrm{id}, \cdot^5\}$.
  Minimality forbids elements of order $\le 4$ in $G \setminus H$, otherwise
  we would be in Case~1 or Case~2. We rule out each remaining $\tau$:
  \begin{itemize}
    \item $\tau = \mathrm{id}$: then $a$ commutes with $H$, and $b := x^{-1}
    a$ satisfies $b^2 = x^{-2} a^2 = x^{-2} \cdot x^2 = e$. Thus $b$ is an
    involution in $G \setminus H$, contradicting minimality.
    \item $\tau = \cdot^5$: then $a x a^{-1} = x^5$, and $(xa)^2 = x \cdot
    a x a^{-1} \cdot a^2 = x \cdot x^5 \cdot x^2 = x^8 = e$. Thus $xa$ is an
    involution in $G \setminus H$, contradicting minimality.
  \end{itemize}
  No surviving extension type arises in Case~3: any group $G$ falling into
  this case is already covered by Cases~1, 2, or 4 via the corresponding
  smaller-order inducing element.

  \smallskip
  \noindent
  \emph{Case 4: $|a| = 16$.} Then $\langle a \rangle = G$, so $G \simeq
  C_{16}$. Note that $a^2$ generates $\langle a \rangle^2 = $ the unique
  subgroup of $G$ of order $8$, which must be $H$. Thus, writing $x = a^2$,
  the extension type associated to $a$ is
  $(C_8, 2, \mathrm{id}, x) = \mathtt{ext\_16\_1}$ (the action is trivial
  because $G$ is abelian). The group $C_{16}$ comes with the realisation
  $\mathtt{realise\_16\_1}$ from
  \texttt{M2rGroup7/SixteenClassification/Blueprints.lean}, and
  Lemma~\ref{lem:realise-of-mulEquiv} applied to a chosen isomorphism
  $G \simeq C_{16}$ transfers this to a realisation of $\mathtt{ext\_16\_1}$
  by $G$.

  \smallskip
  Combining the four cases, $G$ realises one of $\mathtt{ext\_16\_1}$,
  $\mathtt{ext\_16\_5}$, $\mathtt{ext\_16\_6}$, $\mathtt{ext\_16\_7}$,
  $\mathtt{ext\_16\_8}$, or $\mathtt{ext\_16\_9}$, as required.
\end{proof}

\section{Case analysis for normal \texorpdfstring{$K_8 = C_4 \times C_2$}{K8 = C4 x C2}}

\begin{lemma}
  \label{lem:realise-with-normal-K8}
  \lean{OrderSixteen.realise_with_normal_K8}
  \leanfile{M2rGroup7/SixteenClassification/Lemma3.lean}
  \uses{lem:realise-from-normal-index-2, lem:realise-of-mulEquiv}
  Let $G$ be a group of order $16$ with a normal subgroup $H \le G$ such that
  $H \simeq K_8 := C_4 \times C_2$. Then $G$ realises one of the extension
  types $\mathtt{ext\_16\_2}$, $\mathtt{ext\_16\_3}$, $\mathtt{ext\_16\_4}$,
  $\mathtt{ext\_16\_10}$, $\mathtt{ext\_16\_11}$, $\mathtt{ext\_16\_12}$, or
  $\mathtt{ext\_16\_13}$.
\end{lemma}

\begin{proof}
  Write $H = \langle x, y \rangle$ with $|x| = 4$, $|y| = 2$, and $xy = yx$,
  so that $H = \langle x \rangle \times \langle y \rangle \simeq C_4 \times
  C_2$. The hypothesis $|G| = 16$ together with $|H| = 8$ gives $[G : H] = 2$,
  so $H$ is normal of index $2$ and Lagrange yields a partition $G = H \sqcup
  Ha$ for any $a \in G \setminus H$.

  By Lemma~\ref{lem:inducing-element-exists}, $G \setminus H$ is non-empty and
  every $b \in G \setminus H$ satisfies $b^2 \in H$, so $|b|$ is even and
  divides $16$. Hence $|b| \in \{2, 4, 8, 16\}$.

  \smallskip
  \noindent
  \emph{Reduction to $|a| \in \{2, 4\}$.} We may assume $G$ has no element of
  order $8$ (and a fortiori no element of order $16$). Indeed, if some
  $g \in G$ has order $8$ then $\langle g \rangle$ is a subgroup of $G$ of
  index $2$, hence normal, and isomorphic to $C_8$. In that situation the
  theorem
  \texttt{realise\_ext\_type\_if\_not\_iso\_to\_C2\_4}~(Theorem~\ref{thm:realise-ext-type-if-not-iso-to-C2-4})
  dispatches via Lemma~\ref{lem:realise-with-normal-C8} rather than the
  present lemma. Choosing $a$ to be an inducing element of minimal order in
  $G \setminus H$ therefore gives $|a| \in \{2, 4\}$.

  Set $v := a^2 \in H$ and let $\tau := \mathrm{conj}_a \in \mathrm{Aut}(H)$.
  By Lemma~\ref{lem:realise-from-normal-index-2}, $G$ realises the extension
  type $(H, 2, \tau, v)$.

  \smallskip
  \noindent
  \emph{Structure of $\mathrm{Aut}(H)$.} $\mathrm{Aut}(K_8) \simeq D_4$ has
  order $8$. The Lean blueprint uses three named involutions:
  \[
    \psi_3 : (x, y) \mapsto (x^{-1}, y), \qquad
    \psi_5 : (x, y) \mapsto (xy, y), \qquad
    \psi_6 : (x, y) \mapsto (x^{-1}, x^2 y).
  \]
  A short computation shows $\psi_3$ is central in $\mathrm{Aut}(H)$ (it
  commutes with both $\psi_5$ and $\psi_6$), and the kernel of the natural map
  $\mathrm{Aut}(H) \to \mathrm{Aut}(H[2]) \simeq S_3$ (where $H[2] = \{e, x^2,
  y, x^2 y\}$ is the unique Klein four-subgroup) is
  $\{\mathrm{id}, \psi_3, \psi_5, \psi_3 \psi_5\}$. Consequently the
  conjugacy classes of involutions in $\mathrm{Aut}(H)$ are
  \[
    \{\mathrm{id}\}, \quad \{\psi_3\}, \quad \{\psi_5, \psi_3 \psi_5\},
    \quad \{\psi_6, \psi_3 \psi_6\},
  \]
  giving four conjugacy classes of order-$\le 2$ elements. The element $x^2$
  is characteristic in $H$ (it is the unique square of an order-$4$ element),
  so it is fixed by every $\alpha \in \mathrm{Aut}(H)$; the orbit of $y$ under
  $\mathrm{Aut}(H)$ is $\{y, x^2 y\}$ (with $\psi_6$ interchanging them).

  \smallskip
  We split on $|a|$.

  \smallskip
  \noindent
  \emph{Case 1: $|a| = 2$.} Then $v = a^2 = e$, so the validity condition
  $\tau^2 = \mathrm{conj}_v = \mathrm{id}$ just says that $\tau$ is an
  involution (or the identity) in $\mathrm{Aut}(H)$. By the conjugacy-class
  description above, the four
  $\mathrm{Aut}(H)$-conjugacy classes of involutions $\cup \{\mathrm{id}\}$
  give, by Lemma~1(c) (formalised as \texttt{conjugateActEquiv}/
  \texttt{extension\_families\_same\_isoClasses}), four equivalence classes
  of extension types $(H, 2, \tau, e)$ with representatives
  \[
    \tau = \mathrm{id},\quad \tau = \psi_3,\quad \tau = \psi_5,\quad \tau =
    \psi_6,
  \]
  matching exactly the four predefined types
  \[
    \mathtt{ext\_16\_10} = (K_8, 2, \mathrm{id}, e),\quad
    \mathtt{ext\_16\_11} = (K_8, 2, \psi_3, e),
  \]
  \[
    \mathtt{ext\_16\_3} = (K_8, 2, \psi_5, e),\quad
    \mathtt{ext\_16\_13} = (K_8, 2, \psi_6, e),
  \]
  realised respectively by $K_8 \times C_2$, $C_2 \times D_4$,
  $(C_2 \times C_2) \rtimes_{c_4} C_4$, and the semidirect product
  $K_8 \rtimes_{\psi_6} C_2$.

  \smallskip
  \noindent
  \emph{Case 2: $|a| = 4$.} Then $v = a^2$ has order $2$ in $H$, so
  $v \in \{x^2, y, x^2 y\}$. Because $\psi_6$ interchanges $y$ and $x^2 y$,
  Lemma~\ref{lem:realise-of-mulEquiv} together with
  \texttt{conjugateActEquiv} reduces the sub-cases $v = y$ and $v = x^2 y$ to
  each other. We may therefore assume $v \in \{x^2, y\}$.

  \smallskip
  \noindent
  \emph{Sub-case 2a: $v = x^2$.} Since $x^2$ is characteristic, every
  $\alpha \in \mathrm{Aut}(H)$ stabilises $v$, and $\tau^2 = \mathrm{conj}_{x^2}
  = \mathrm{id}$ (as $H$ is abelian) forces $\tau$ to be the identity or an
  involution. By Lemma~1(c) we may take $\tau$ to be a representative of one
  of the four conjugacy classes $\{\mathrm{id}\}, \{\psi_3\}, \{\psi_5\},
  \{\psi_6\}$. For each $b = x^k y^l a \in G \setminus H$ a direct computation
  gives
  \[
    (x^k y^l a)^2 = x^k y^l \cdot \tau(x^k y^l) \cdot v.
  \]
  \begin{itemize}
    \item $\tau = \mathrm{id}$: $\tau(x^k y^l) = x^k y^l$, so $(b)^2 = x^{2k}
    y^{2l} \cdot x^2 = x^{2k + 2}$. Taking $k = -1, l = 0$ gives $(x^{-1} a)^2
    = e$, so $x^{-1} a$ is an involution in $G \setminus H$, contradicting
    minimality.
    \item $\tau = \psi_3$: $\tau(x^k y^l) = x^{-k} y^l$, so $b \cdot \tau(b) =
    y^{2l} = e$ and hence $(b)^2 = x^2$ for every $h = x^k y^l$. Thus every
    element of $G \setminus H$ has order $4$ (since $(b)^4 = x^4 = e$),
    confirming minimality. The resulting extension type is
    $\mathtt{ext\_16\_12} = (K_8, 2, \psi_3, x^2)$, realised by
    $C_2 \times Q_2$.
    \item $\tau = \psi_5$: $\tau(x^k y^l) = x^k y^{k+l}$, so $b \cdot \tau(b)
    = x^{2k} y^{2l + k} = x^{2k} y^{k \bmod 2}$, giving
    $(b)^2 = x^{2k + 2} y^{k \bmod 2}$. This equals $e$ iff $2k + 2 \equiv 0
    \pmod 4$ (i.e.\ $k$ odd) and $k \equiv 0 \pmod 2$ (i.e.\ $k$ even),
    which is impossible. Hence every element of $G \setminus H$ has order
    exactly $4$ and minimality is preserved. The resulting extension type is
    $\mathtt{ext\_16\_4} = (K_8, 2, \psi_5, x^2)$, realised by $C_4 \rtimes
    _{\mathrm{inv}} C_4$.
    \item $\tau = \psi_6$: $\tau(x^k y^l) = x^{-k} (x^2 y)^l = x^{2l - k}
    y^l$, so $b \cdot \tau(b) = x^{2l} y^{2l} = x^{2l}$, giving $(b)^2 =
    x^{2l + 2}$. Taking $l = 1, k = 0$ gives $(y a)^2 = e$, so $y a$ is an
    involution in $G \setminus H$, contradicting minimality.
  \end{itemize}
  Hence Sub-case~2a yields the two surviving extension types
  $\mathtt{ext\_16\_4}$ and $\mathtt{ext\_16\_12}$.

  \smallskip
  \noindent
  \emph{Sub-case 2b: $v = y$.} The stabiliser of $y$ in $\mathrm{Aut}(H)$ is
  the kernel of the map $\mathrm{Aut}(H) \to \mathrm{Aut}(H[2])$ described
  above, namely $V = \{\mathrm{id}, \psi_3, \psi_5, \psi_3 \psi_5\}$. Hence
  validity $\tau(v) = v$ together with $\tau^2 = \mathrm{conj}_y = \mathrm{id}$
  restricts $\tau$ to $V$. Lemma~1(c) only allows conjugation by elements of
  $V$, but $V$ is abelian, so each of the four $\tau \in V$ must be analysed
  separately. As in Sub-case~2a we use
  \[
    (x^k y^l a)^2 = x^k y^l \cdot \tau(x^k y^l) \cdot v.
  \]
  \begin{itemize}
    \item $\tau = \mathrm{id}$: $b \cdot \tau(b) = x^{2k} y^{2l} = x^{2k}$, so
    $(b)^2 = x^{2k} y$. For $k$ even this is $y$, and for $k$ odd this is
    $x^2 y$. In all cases $(b)^2 \neq e$, so every element of $G \setminus H$
    has order $4$ and minimality is preserved. The resulting extension type
    is $\mathtt{ext\_16\_2} = (K_8, 2, \mathrm{id}, y)$, realised by $C_4
    \times C_4$.
    \item $\tau = \psi_5$: by Lemma~1(b) (the cocycle-shift
    \texttt{extension\_families\_same\_isoClasses}), $(\psi_5, y) \sim (\psi_5,
    y \cdot h \cdot \psi_5(h))$ for any $h \in H$. Taking $h = x$ gives
    $y \cdot x \cdot \psi_5(x) = y \cdot x \cdot xy = x^2$, so $(\psi_5, y) \sim
    (\psi_5, x^2)$, which is $\mathtt{ext\_16\_4}$ and is already covered by
    Sub-case~2a.
    \item $\tau = \psi_3 \psi_5$: write $\sigma = \psi_3 \psi_5$, so
    $\sigma(x^k y^l) = x^{-k} y^{k + l}$ and $b \cdot \sigma(b) = y^{2l + k} =
    y^{k \bmod 2}$, giving $(b)^2 = y \cdot y^{k \bmod 2}$. Taking $k = 1$
    gives $(x a)^2 = e$, so $x a$ is an involution in $G \setminus H$,
    contradicting minimality.
    \item $\tau = \psi_3$: $b \cdot \tau(b) = y^{2l} = e$, so $(b)^2 = y$ and
    every element of $G \setminus H$ has order $4$. Minimality is preserved,
    but the resulting extension type $(K_8, 2, \psi_3, y)$ is not one of the
    seven $K_8$-based predefined types. However, the group $G$ defined by this
    extension data is presented by
    \[
      \langle X, a \mid X^4 = a^4 = e,\ a X a^{-1} = X^{-1} \rangle,
    \]
    after eliminating the redundant generator $y = a^2$ (which lies in
    $Z(G)$ since $a^2$ commutes with $a$ and with $X$, the latter because
    $a^2 X a^{-2} = a X^{-1} a^{-1} = X$). This presentation is exactly
    that of the realising group $C_4 \rtimes_{\mathrm{inv}} C_4$ of
    $\mathtt{ext\_16\_4}$. Lemma~\ref{lem:realise-of-mulEquiv} applied to the
    induced isomorphism $G \simeq C_4 \rtimes_{\mathrm{inv}} C_4$ transfers
    $\mathtt{realise\_16\_4}$ to $G$, so this sub-sub-case also falls under
    $\mathtt{ext\_16\_4}$.
  \end{itemize}
  Hence Sub-case~2b yields the surviving new extension type
  $\mathtt{ext\_16\_2}$ (with $\tau = \mathrm{id}$), while the remaining
  $\tau \in V$ are either absorbed into $\mathtt{ext\_16\_4}$ or killed by
  minimality.

  \smallskip
  Combining Cases~1 and~2, $G$ realises one of
  \[
    \mathtt{ext\_16\_2},\ \mathtt{ext\_16\_3},\ \mathtt{ext\_16\_4},\
    \mathtt{ext\_16\_10},\ \mathtt{ext\_16\_11},\ \mathtt{ext\_16\_12},\
    \mathtt{ext\_16\_13},
  \]
  as required.
\end{proof}

\section{Main lemma (Wild's Lemma 3)}

\begin{theorem}
  \label{thm:realise-ext-type-if-not-iso-to-C2-4}
  \lean{OrderSixteen.realise_ext_type_if_not_iso_to_C2_4}
  \leanfile{M2rGroup7/SixteenClassification/Classification.lean}
  \uses{lem:realise-with-normal-C8, lem:realise-with-normal-K8}
  Let $G$ be a group of order $16$ that is not isomorphic to $(C_2)^4$. Then
  $G$ realises one of the extension types $\mathtt{ext\_16\_1}, \dots,
  \mathtt{ext\_16\_13}$.
\end{theorem}

\begin{proof}
  By the Lean lemma \texttt{exists\_normal\_C8\_or\_C4\_C2}, $G$ has either a
  normal subgroup isomorphic to $C_8$ or a normal subgroup isomorphic to
  $K_8 = C_4 \times C_2$. The two cases are handled respectively by
  Lemma~\ref{lem:realise-with-normal-C8} and
  Lemma~\ref{lem:realise-with-normal-K8}, exhausting the possibilities and
  yielding one of the thirteen extension types.
\end{proof}

\end{document}


%%FILE numina/blueprints/4q-classification-refactor/4q-classification-refactor.tex
% Blueprint: 4q classification refactor
% Created by Numina

\begin{document}

\section{Setup}

We classify finite groups $G$ of order $4q$, where $q$ is a prime with $q \geq
3$. Equivalently, $|G| = p^2 q$ for $p = 2$. The full classification produces
five generic isomorphism types and one extra type when $q = 3$:
\[
\begin{aligned}
& C_{4q}, \quad C_2 \times C_2 \times C_q, \\
& C_q \rtimes_{\varphi_1} C_4, \quad
  C_q \rtimes_{\varphi_2} C_4 \ (\text{only when } q \equiv 1 \pmod 4), \\
& C_q \rtimes (C_2 \times C_2), \quad
  (C_2 \times C_2) \rtimes C_3 \ (\text{only when } q = 3).
\end{aligned}
\]
The actions $\varphi_1, \varphi_2$ are the canonical actions of $C_4$ on $C_q$
with image of order $2$ and $4$ respectively; the action of $C_2 \times C_2$ on
$C_q$ is canonical with image of order $2$, giving the dihedral group
$D_{2q}$; the action of $C_3$ on $C_2 \times C_2$ is the canonical order-$3$
permutation (giving $A_4$ when $q = 3$).

This refactor moves the heavy lifting into a single classification engine
\texttt{classify\_sdp} that operates on \emph{abstract} cyclic groups and uses
a canonical action transported to those abstract groups
(\texttt{sdpCanonicalAction}). At call sites, $\varphi$ is conjugated into the
concrete cyclic groups \emph{before} classification, so the output isomorphism
already lives in the concrete cyclic groups by definitional equality. This
eliminates the long chain of transport bridges previously needed.

\section{Sylow structure for $p^2 q$ groups}

\begin{lemma}[A group of order $p^2 q$ has a normal Sylow subgroup]
    \label{lem:p2q-normal-sylow}
    Let $p, q$ be distinct primes and let $G$ be a finite group of order
    $p^2 q$. Then either the number of Sylow $p$-subgroups of $G$ is $1$ (so
    the unique Sylow $p$-subgroup is normal) or the number of Sylow
    $q$-subgroups of $G$ is $1$ (so the unique Sylow $q$-subgroup is normal).
\end{lemma}
\begin{proof}
    Standard Sylow counting. Already established as
    \texttt{p2q\_group\_has\_normal\_sylow\_subgroup} in
    \texttt{M2rGroup7.Lemmas.SylowUtils}.
\end{proof}

\begin{lemma}[Sylow cardinality and index]
    \label{lem:sylow-card-index}
    Let $p, q$ be distinct primes, $a, b \in \mathbb{N}$, and let $G$ be a
    finite group with $|G| = p^a q^b$. Then every Sylow $p$-subgroup of $G$
    has order $p^a$ and index $q^b$.
\end{lemma}
\begin{proof}
    From $|G| = p^a q^b$ and coprimality of $p^a$ and $q^b$, the multiplicity
    of $p$ in $|G|$ is $a$, so every Sylow $p$-subgroup has order $p^a$ by
    Mathlib's \texttt{Sylow.card\_eq\_multiplicity}. The index is then
    $|G|/p^a = q^b$ by Lagrange. Already established as
    \texttt{sylow\_card\_eq} and \texttt{sylow\_index\_eq}.
\end{proof}

\section{Background results}

\begin{lemma}[Cyclic classification of prime order]
    \label{lem:prime-classification}
    If $|G| = n$ with $n$ prime, then $G \cong C_n$.
\end{lemma}
\begin{proof}
    A group of prime order is cyclic. Two cyclic groups of equal finite order
    are isomorphic via Mathlib's \texttt{mulEquivOfCyclicCardEq}. Already
    established as \texttt{prime\_classification} in
    \texttt{M2rGroup7.Classification}.
\end{proof}

\begin{lemma}[Classification of groups of order $p^2$]
    \label{lem:p-squared-classification}
    For any prime $p$, a group of order $p^2$ is isomorphic to either
    $C_{p^2}$ or $C_p \times C_p$.
\end{lemma}
\begin{proof}
    Any group of order $p^2$ is abelian; either it is cyclic (and isomorphic
    to $C_{p^2}$) or it has exponent $p$, in which case it is a $2$-dimensional
    $\mathbb{F}_p$-vector space. Already established as
    \texttt{p\_squared\_classification}.
\end{proof}

\begin{lemma}[Aut$(C_{p^2}) \cong C_{p(p-1)}$]
    \label{lem:aut-cyclic-p2}
    For any prime $p$, $\mathrm{MulAut}(C_{p^2}) \cong C_{p(p-1)}$.
\end{lemma}
\begin{proof}
    $\mathrm{MulAut}(C_{p^2}) \cong (\mathbb{Z}/p^2 \mathbb{Z})^\times$, which is
    cyclic of order $\varphi(p^2) = p(p-1)$. Already established as
    \texttt{aut\_of\_cyclic\_p2}.
\end{proof}

\begin{lemma}[Aut$(C_p \times C_p) \cong \mathrm{GL}_2(\mathbb{F}_p)$]
    \label{lem:aut-CpCp}
    For any prime $p$, $\mathrm{MulAut}(C_p \times C_p) \cong
    \mathrm{GL}_2(\mathbb{F}_p)$.
\end{lemma}
\begin{proof}
    $C_p \times C_p$ is a $2$-dimensional $\mathbb{F}_p$-vector space and
    every additive automorphism is automatically $\mathbb{F}_p$-linear.
    Already established as \texttt{aut\_of\_CpCp}.
\end{proof}

\begin{lemma}[$\mathrm{GL}_2(\mathbb{F}_2) \cong D_3$]
    \label{lem:GL2F2-isoS3}
    $\mathrm{GL}_2(\mathbb{F}_2)$ is isomorphic to the dihedral group of order
    $6$.
\end{lemma}
\begin{proof}
    A direct computation: $|\mathrm{GL}_2(\mathbb{F}_2)| = 6$ and
    $\mathrm{GL}_2(\mathbb{F}_2)$ is non-abelian. Already established as
    \texttt{GL2F2\_isoS3}.
\end{proof}

\begin{lemma}[Trivial homomorphism from coprime cardinalities]
    \label{lem:eq-one-of-coprime-card}
    Let $G, H$ be finite groups with $\gcd(|G|, |H|) = 1$. Then every group
    homomorphism $\varphi : G \to H$ satisfies $\varphi = 1$.
\end{lemma}
\begin{proof}
    The image $\varphi(G)$ is a subgroup of $H$ whose cardinality divides
    $\gcd(|G|,|H|) = 1$, so $|\varphi(G)| = 1$ and $\varphi$ is trivial.
    Already established as \texttt{eq\_one\_of\_coprime\_card}.
\end{proof}

\begin{lemma}[Semidirect product with trivial action is a direct product]
    \label{lem:sdp-trivial-action}
    Let $N, K$ be groups and let $\varphi : K \to \mathrm{Aut}(N)$ with
    $\varphi = 1$. Then $N \rtimes_\varphi K \cong N \times K$.
\end{lemma}
\begin{proof}
    The multiplication on $N \rtimes_\varphi K$,
    $(n_1, k_1)(n_2, k_2) = (n_1 \cdot \varphi(k_1)(n_2),\ k_1 k_2)$,
    reduces to $(n_1 n_2, k_1 k_2)$. Already established as
    \texttt{SemidirectProduct.mulEquivOfTrivialAction}.
\end{proof}

\begin{lemma}[Coprime cyclic direct product]
    \label{lem:cyclic-prod}
    For positive integers $m, n$ with $\gcd(m, n) = 1$,
    $C_m \times C_n \cong C_{mn}$.
\end{lemma}
\begin{proof}
    Already established as \texttt{CyclicGroup.prodMulEquiv}.
\end{proof}

\begin{lemma}[Schur--Zassenhaus complement]
    \label{lem:schur-zassenhaus}
    Let $G$ be a finite group and $N \trianglelefteq G$ a normal subgroup
    with $\gcd(|N|, [G:N]) = 1$. Then there exists a subgroup $K \leq G$
    with $N \cap K = \{1\}$ and $NK = G$.
\end{lemma}
\begin{proof}
    Mathlib's \texttt{Subgroup.exists\_right\_complement'\_of\_coprime}.
\end{proof}

\begin{lemma}[Semidirect product from a complement]
    \label{lem:sdp-from-complement}
    Let $G$ be a group with $N \trianglelefteq G$ and let $K \leq G$ be a
    complement of $N$. Then $G \cong N \rtimes_\varphi K$ where $\varphi$ is
    the conjugation action of $K$ on $N$.
\end{lemma}
\begin{proof}
    Mathlib's \texttt{SemidirectProduct.mulEquivSubgroup}.
\end{proof}

\begin{lemma}[Transport of a semidirect product]
    \label{lem:sdp-congr}
    Let $f_N : N_1 \cong N_2$ and $f_K : K_1 \cong K_2$ be group isomorphisms
    and let $\varphi_1 : K_1 \to \mathrm{Aut}(N_1)$. Set
    $\varphi_2 := (\mathrm{MulAut.congr}\,f_N) \circ \varphi_1 \circ f_K^{-1}$.
    Then $N_1 \rtimes_{\varphi_1} K_1 \cong N_2 \rtimes_{\varphi_2} K_2$.
\end{lemma}
\begin{proof}
    \texttt{SemidirectProduct.congr'}; the underlying map sends
    $(n, k)$ to $(f_N(n), f_K(k))$.
\end{proof}

\section{Classification engine on cyclic groups}

The refactor introduces a single classification engine
\texttt{classify\_sdp} that takes abstract cyclic groups $N$ and $K$ with
$|N| = q^n$ and $|K| = p^m$ and any $\varphi : K \to \mathrm{Aut}(N)$, and
produces a unique invariant $r$ together with an isomorphism to a canonical
semidirect product. The canonical action is built directly in the abstract
groups so that no transport bridges are needed at call sites.

\begin{theorem}[Classification of semidirect products of cyclic groups]
    \label{thm:classify-sdp}
    \uses{lem:sdp-congr}
    Let $p, q$ be distinct primes with $q$ odd, let $m \geq 1$ and $n \geq 1$,
    and let $N$, $K$ be (abstract) cyclic groups with $|N| = q^n$ and
    $|K| = p^m$. For any homomorphism $\varphi : K \to \mathrm{Aut}(N)$ there
    exists a unique $r \in \{0, 1, \dots, \min(m, v_p(q-1))\}$ together with an
    isomorphism
    \[
        N \rtimes_\varphi K \ \cong\ N \rtimes_{\psi_r} K,
    \]
    where $\psi_r := \mathtt{sdpCanonicalAction}(\dots, r) : K \to
    \mathrm{Aut}(N)$ is the canonical action of image-order $p^r$.
\end{theorem}
\begin{proof}
    A two-step bridge using Lemma~\ref{lem:sdp-congr}: conjugate $\varphi$
    across the unique iso $N \cong C_{q^n}$ and $K \cong C_{p^m}$, apply the
    concrete classification on $C_{q^n} \rtimes C_{p^m}$, and conjugate back.
    Implemented as \texttt{classify\_sdp} in
    \texttt{M2rGroup7.P2qClassification.CycPGroupClassification}.
\end{proof}

\begin{lemma}[Range cardinality of $\psi_r$]
    \label{lem:sdp-canonical-range}
    \uses{thm:classify-sdp}
    The image $\psi_r(K)$ has cardinality exactly $p^r$.
\end{lemma}
\begin{proof}
    Already established as \texttt{sdpCanonicalAction\_range\_card}.
\end{proof}

\begin{lemma}[Iso of two homomorphisms with equal images]
    \label{lem:sdp-iso-range-eq}
    Let $p$ be prime, $N$ be a finite group with $\gcd(|N|, p) = 1$, $K$ be a
    cyclic $p$-group with $|K| = p$, and let
    $f_1, f_2 : K \to \mathrm{Aut}(N)$ have the same image. Then
    $N \rtimes_{f_1} K \cong N \rtimes_{f_2} K$.
\end{lemma}
\begin{proof}
    Already established as \texttt{semidirectProduct\_iso\_if\_range\_eq}.
\end{proof}

\section{Canonical actions for $4q$ groups}

The three canonical actions used in the $4q$ classification are obtained by
specialising \texttt{sdpCanonicalAction} to the concrete cyclic groups
$C_q$ and $C_4$ (or $C_2$).

\begin{definition}[$\varphi_1 : C_4 \to \mathrm{Aut}(C_q)$, image of order $2$]
    \label{def:canonicalC4OnCq-r1}
    \uses{thm:classify-sdp}
    For an odd prime $q$, define
    \[
        \mathtt{canonicalC4OnCqAction} :=
            \mathtt{sdpCanonicalAction}\bigl(\dots, m = 2,\, n = 1,\, r = 1\bigr)
            \,:\, C_4 \to \mathrm{Aut}(C_q).
    \]
\end{definition}

\begin{definition}[$\varphi_2 : C_4 \to \mathrm{Aut}(C_q)$, image of order $4$]
    \label{def:canonicalC4OnCq-r2}
    \uses{thm:classify-sdp}
    For an odd prime $q$ with $q \equiv 1 \pmod 4$, define
    \[
        \mathtt{canonicalC4OnCqAction\_r2} :=
            \mathtt{sdpCanonicalAction}\bigl(\dots, m = 2,\, n = 1,\, r = 2\bigr).
    \]
    The hypothesis $r = 2 \leq \min(2, v_2(q - 1))$ requires $v_2(q-1) \geq 2$,
    equivalently $q \equiv 1 \pmod 4$.
\end{definition}

\begin{definition}[$\varphi : C_2 \times C_2 \to \mathrm{Aut}(C_q)$, image of order $2$]
    \label{def:canonicalC2C2OnCq}
    \uses{thm:classify-sdp}
    For an odd prime $q$, define
    \[
        \mathtt{canonicalC2C2OnCqAction} :=
            \mathtt{sdpCanonicalAction}\bigl(\dots, m = 1,\, n = 1,\, r = 1\bigr)
            \circ \mathrm{fst},
    \]
    where $\mathrm{fst} : C_2 \times C_2 \to C_2$ is projection on the first
    factor. The resulting semidirect product $C_q \rtimes (C_2 \times C_2)$
    is isomorphic to the dihedral group $D_{2q}$.
\end{definition}

\section{The order-$3$ canonical action (case $q = 3$)}

When $q = 3$ and the Sylow $2$-subgroup is $C_2 \times C_2$, the action
$C_3 \to \mathrm{Aut}(C_2 \times C_2) \cong S_3$ can be nontrivial because
$S_3$ has elements of order $3$. The resulting semidirect product is $A_4$.

\begin{definition}[Order-$3$ automorphism of $C_2 \times C_2$]
    \label{def:c2c2OrderThreeAut}
    Let $\sigma : C_2 \times C_2 \to C_2 \times C_2$ be the automorphism
    $(x, y) \mapsto (xy, x)$. Then $\sigma^3 = \mathrm{id}$ and $\sigma$ has
    order exactly $3$ in $\mathrm{MulAut}(C_2 \times C_2)$.
\end{definition}
\begin{lemma}
    \label{lem:orderOf-c2c2OrderThreeAut}
    \uses{def:c2c2OrderThreeAut}
    $\mathrm{ord}(\sigma) = 3$.
\end{lemma}
\begin{proof}
    Verify $\sigma^3 = \mathrm{id}$ by direct computation in $C_2 \times C_2$
    using that every element has order dividing $2$, then exhibit a witness
    $(g, 1)$ with $g$ a generator of $C_2$ where $\sigma(g, 1) = (g, g)
    \neq (g, 1)$ to show $\sigma \neq \mathrm{id}$.
\end{proof}

\begin{definition}[$\psi : C_3 \to \mathrm{Aut}(C_2 \times C_2)$ sending generator to $\sigma$]
    \label{def:canonicalC3OnC2C2}
    \uses{lem:orderOf-c2c2OrderThreeAut}
    Define $\mathtt{canonicalC3OnC2C2Action} : C_3 \to
    \mathrm{MulAut}(C_2 \times C_2)$ by extending the assignment
    (generator of $C_3$) $\mapsto \sigma$ to a group homomorphism, which is
    well-defined because $|\sigma| = 3 = |C_3|$.
\end{definition}

\begin{lemma}[Range cardinality of $\psi$]
    \label{lem:canonicalC3OnC2C2-range}
    \uses{def:canonicalC3OnC2C2, lem:orderOf-c2c2OrderThreeAut}
    The image of $\mathtt{canonicalC3OnC2C2Action}$ has cardinality $3$.
\end{lemma}
\begin{proof}
    The range equals $\langle \sigma \rangle$, whose cardinality is
    $\mathrm{ord}(\sigma) = 3$.
\end{proof}

\begin{lemma}[Uniqueness of semidirect product $C_3 \to \mathrm{Aut}(C_2 \times C_2)$ with image of order $3$]
    \label{lem:sdp-C3-on-C2C2-iso}
    \uses{lem:canonicalC3OnC2C2-range}
    Let $f_1, f_2 : C_3 \to \mathrm{Aut}(C_2 \times C_2)$ both have image of
    cardinality $3$. Then $(C_2 \times C_2) \rtimes_{f_1} C_3 \cong
    (C_2 \times C_2) \rtimes_{f_2} C_3$.
\end{lemma}
\begin{proof}
    $\mathrm{MulAut}(C_2 \times C_2) \cong D_3$ has a unique Sylow
    $3$-subgroup (since $|D_3| = 6 = 3 \cdot 2$ and $n_3 \equiv 1
    \pmod 3$ with $n_3 \mid 2$ forces $n_3 = 1$). Hence both images coincide
    with this Sylow $3$-subgroup. Apply
    Lemma~\ref{lem:sdp-iso-range-eq} with $p = 3$.
\end{proof}

\section{The main classification}

\begin{theorem}[Classification of groups of order $4q$ for $q \geq 3$ prime]
    \label{thm:classification-4q}
    \uses{lem:p2q-normal-sylow, lem:sylow-card-index,
          lem:prime-classification, lem:p-squared-classification,
          lem:aut-cyclic-p2, lem:aut-CpCp, lem:GL2F2-isoS3,
          lem:eq-one-of-coprime-card, lem:sdp-trivial-action,
          lem:cyclic-prod, lem:schur-zassenhaus, lem:sdp-from-complement,
          lem:sdp-congr,
          thm:classify-sdp, lem:sdp-canonical-range,
          def:canonicalC4OnCq-r1, def:canonicalC4OnCq-r2,
          def:canonicalC2C2OnCq, def:canonicalC3OnC2C2,
          lem:sdp-C3-on-C2C2-iso}
    Let $q$ be a prime with $q \geq 3$ and let $G$ be a group with $|G| = 4q$.
    Then exactly one of the following holds:
    \begin{enumerate}
        \item $G \cong C_{4q}$.
        \item $G \cong C_2 \times C_2 \times C_q$.
        \item $G \cong C_q \rtimes_{\varphi_1} C_4$, where $\varphi_1$ is
            $\mathtt{canonicalC4OnCqAction}$ (image of order $2$).
        \item $q \equiv 1 \pmod 4$ and $G \cong C_q \rtimes_{\varphi_2} C_4$,
            where $\varphi_2$ is $\mathtt{canonicalC4OnCqAction\_r2}$ (image
            of order $4$).
        \item $G \cong C_q \rtimes_\varphi (C_2 \times C_2)$, where $\varphi$
            is $\mathtt{canonicalC2C2OnCqAction}$ (image of order $2$);
            this group is isomorphic to $D_{2q}$.
        \item $q = 3$ and $G \cong (C_2 \times C_2) \rtimes_\psi C_3$, where
            $\psi$ is $\mathtt{canonicalC3OnC2C2Action}$; this group is
            isomorphic to $A_4$.
    \end{enumerate}
    Implemented as \texttt{classification\_4q} in
    \texttt{M2rGroup7.P2qClassification.FourQClassification}.
\end{theorem}
\begin{proof}
    \uses{lem:p2q-normal-sylow, lem:sylow-card-index,
          lem:prime-classification, lem:p-squared-classification,
          lem:aut-cyclic-p2, lem:aut-CpCp, lem:GL2F2-isoS3,
          lem:eq-one-of-coprime-card, lem:sdp-trivial-action,
          lem:cyclic-prod, lem:schur-zassenhaus, lem:sdp-from-complement,
          lem:sdp-congr,
          thm:classify-sdp, lem:sdp-canonical-range,
          def:canonicalC4OnCq-r1, def:canonicalC4OnCq-r2,
          def:canonicalC2C2OnCq, def:canonicalC3OnC2C2,
          lem:sdp-C3-on-C2C2-iso}
    By Lemma~\ref{lem:p2q-normal-sylow} applied with $p = 2$, either the
    Sylow $2$-subgroup or the Sylow $q$-subgroup of $G$ is normal.

    \emph{Case 1: normal Sylow $2$-subgroup $P$.} By
    Lemma~\ref{lem:sylow-card-index}, $|P| = 4$ and $[G:P] = q$. By
    Lemma~\ref{lem:schur-zassenhaus} (with $\gcd(4, q) = 1$) there is a
    complement $K \leq G$ with $|K| = q$. By
    Lemma~\ref{lem:sdp-from-complement}, $G \cong P \rtimes_\varphi K$ for
    some $\varphi$. By Lemma~\ref{lem:prime-classification}, $K \cong C_q$.

    \begin{itemize}
        \item If $P \cong C_4$: by Lemma~\ref{lem:aut-cyclic-p2},
            $\mathrm{Aut}(P) \cong C_2$. Since $\gcd(q, 2) = 1$, by
            Lemma~\ref{lem:eq-one-of-coprime-card}, $\varphi = 1$. Combine
            Lemma~\ref{lem:sdp-trivial-action} and
            Lemma~\ref{lem:cyclic-prod} to get $G \cong C_{4q}$ (outcome 1).

        \item If $P \cong C_2 \times C_2$: by Lemma~\ref{lem:aut-CpCp} and
            Lemma~\ref{lem:GL2F2-isoS3}, $|\mathrm{Aut}(P)| = 6$.

            If $q > 3$, then $\gcd(q, 6) = 1$ and
            Lemma~\ref{lem:eq-one-of-coprime-card} forces $\varphi = 1$,
            giving $G \cong C_2 \times C_2 \times C_q$ (outcome 2).

            If $q = 3$, we split on whether $\varphi$ is trivial.
            If yes, $G \cong C_2 \times C_2 \times C_3$ (outcome 2). If no,
            transport $\varphi$ along $P \cong C_2 \times C_2$ and
            $K \cong C_3$ to a nontrivial
            $\varphi' : C_3 \to \mathrm{Aut}(C_2 \times C_2)$. Its image has
            cardinality $3$ (since the image is a non-trivial subgroup of the
            cyclic $C_3$). By Lemma~\ref{lem:canonicalC3OnC2C2-range} and
            Lemma~\ref{lem:sdp-C3-on-C2C2-iso}, $G$ is isomorphic to the
            semidirect product with $\mathtt{canonicalC3OnC2C2Action}$
            (outcome 6).
    \end{itemize}

    \emph{Case 2: normal Sylow $q$-subgroup $Q$.} By
    Lemma~\ref{lem:sylow-card-index}, $|Q| = q$ and $[G:Q] = 4$. By
    Lemma~\ref{lem:schur-zassenhaus} there is a complement $K$ with
    $|K| = 4$, and by Lemma~\ref{lem:sdp-from-complement},
    $G \cong Q \rtimes_\varphi K$ with $\varphi : K \to \mathrm{Aut}(Q)$. By
    Lemma~\ref{lem:prime-classification}, $Q \cong C_q$.

    \begin{itemize}
        \item If $K \cong C_4$: transport $\varphi$ across $Q \cong C_q$
            and $K \cong C_4$ (Lemma~\ref{lem:sdp-congr}) to obtain
            $\varphi' : C_4 \to \mathrm{Aut}(C_q)$. Apply
            Theorem~\ref{thm:classify-sdp} with $p = 2$, $m = 2$, $n = 1$ to
            get a unique
            $r \in \{0, 1, 2\} \cap \{0, \dots, \min(2, v_2(q-1))\}$ and an
            isomorphism to the canonical $\psi_r$-semidirect product.
            \begin{itemize}
                \item $r = 0$: $\psi_0$ has image of cardinality $1$ by
                    Lemma~\ref{lem:sdp-canonical-range}, so $\psi_0 = 1$;
                    Lemma~\ref{lem:sdp-trivial-action} and
                    Lemma~\ref{lem:cyclic-prod} give
                    $G \cong C_q \times C_4 \cong C_{4q}$ (outcome 1).
                \item $r = 1$: $\psi_1$ equals
                    $\mathtt{canonicalC4OnCqAction}$ by definitional equality
                    (same head function, same explicit arguments, proof
                    arguments equal by proof irrelevance), giving outcome 3.
                \item $r = 2$: forces $\min(2, v_2(q-1)) \geq 2$ hence
                    $v_2(q-1) \geq 2$ hence $q \equiv 1 \pmod 4$; $\psi_2$
                    equals $\mathtt{canonicalC4OnCqAction\_r2}$ by
                    definitional equality, giving outcome 4.
            \end{itemize}

        \item If $K \cong C_2 \times C_2$: if $\varphi = 1$ then
            Lemma~\ref{lem:sdp-trivial-action} gives
            $G \cong C_q \times C_2 \times C_2 \cong C_2 \times C_2 \times C_q$
            (outcome 2). Otherwise, transport $\varphi$ to
            $\varphi' : C_2 \times C_2 \to \mathrm{Aut}(C_q)$; its image is a
            subgroup of $\mathrm{Aut}(C_q) \cong C_{q-1}$ and thus cyclic.
            The image cardinality divides $|C_2 \times C_2| = 4$ and divides
            $|\mathrm{Aut}(C_q)| = q - 1$, but also (as a hom from
            $C_2 \times C_2$ into a cyclic group) divides $2$; combined with
            $\varphi' \neq 1$, the image has cardinality exactly $2$.
            $\mathtt{canonicalC2C2OnCqAction}$ has the same image cardinality
            (Definition~\ref{def:canonicalC2C2OnCq}), so by
            \texttt{semidirectProduct\_CpCp\_iso} (a $C_p \times C_p$ analogue
            of Lemma~\ref{lem:sdp-iso-range-eq}; see
            \texttt{M2rGroup7.Lemmas.ClassificationUtils}) the two semidirect
            products are isomorphic, giving outcome 5.
    \end{itemize}
\end{proof}

\section{Notes on the refactor}

The chief design choice in this refactor is to define
\texttt{canonicalC4OnCqAction}, \texttt{canonicalC4OnCqAction\_r2}, and
\texttt{canonicalC2C2OnCqAction} as direct applications of
\texttt{sdpCanonicalAction} to the concrete cyclic groups $C_q$ and $C_4$ (or
$C_2$). All three are marked \texttt{@[reducible]} so that the elaborator can
unify, by congruence and proof irrelevance, the term
\texttt{sdpCanonicalAction \dots r \dots} produced by
\texttt{classify\_sdp} with each canonical action at the appropriate value of
$r$.

In the Case 2 / $K \cong C_4$ branch the proof destructures the output
$r : \mathrm{Fin}(\min(2, v_2(q-1)) + 1)$ as $\langle r_{\mathrm{val}},
h_{\mathrm{lt}} \rangle$ and applies \texttt{interval\_cases r\_val}, so
that within each case $r$ in the type of the witness iso is literally
substituted by $0$, $1$, or $2$. This eliminates the need for an explicit
rewrite chain and allows the $r = 1$ and $r = 2$ cases to close with
$\mathtt{h\_iso\_g\_q\_k.symm.trans\;(h\_bridge.trans\;e)}$.

\end{document}


%%FILE numina/blueprints/4q-classification-refactor/source/4q-classification-refactor-source.tex
Mostly fixing errors in 4qClassification.lean

The intended refactor (pasting from chat with Claude):

"Yes, that's exactly right, and you should feel confident about it. Let me state it precisely to make sure we're aligned:

  You have φ : ↥K →* MulAut ↥Q (abstract Sylow subgroups). You want the theorem output to mention CyclicGroup q and CyclicGroup 4. So before calling classify_sdp, you use eQ : ↥Q ≃* CyclicGroup q and
  h_K_C4.some : ↥K ≃* CyclicGroup 4 to conjugate φ into φ' : CyclicGroup 4 →* MulAut (CyclicGroup q), and get the bridge iso ↥Q ⋊ φ ≃* CyclicGroup q ⋊ φ'.

  Now classify_sdp is fed φ' operating on the concrete types (CyclicGroup q, CyclicGroup 4). The output says CyclicGroup q ⋊ φ' ≃* CyclicGroup q ⋊ (sdpCanonicalAction r for CyclicGroup q, CyclicGroup 4).

  On the "exactly equal vs extremely close" question: if canonicalC4OnCqAction is defined as sdpCanonicalAction applied to (CyclicGroup q, CyclicGroup 4) with exactly the same hN and hK that classify_sdp
  uses, they are definitionally equal by congruence with no unfolding — same function, same explicit arguments. So it's "exactly equal" from Lean's perspective, and the r=1 proof step is literally just
  chaining the three isos together.

  The key insight you had (that I missed): the transport step belongs before classify_sdp, not after. That eliminates the need for sdpCanonicalAction_transport entirely."

"Yes, that's exactly what I mean. canonicalC4OnCqAction should literally be defined as sdpCanonicalAction instantiated at the concrete cyclic groups:

  noncomputable def canonicalC4OnCqAction {q : ℕ} [hq : Fact q.Prime] (h : q ≠ 2) :
      CyclicGroup 4 →* MulAut (CyclicGroup q) :=
    haveI : Fact (Nat.Prime 2) := ⟨by norm_num⟩
    sdpCanonicalAction (show 2 ≠ q from by omega) h 2 1 Nat.one_pos
      (by rw [card_cyclicGroup, pow_one])   -- Nat.card (CyclicGroup q) = q^1
      (by rw [card_cyclicGroup]; norm_num)  -- Nat.card (CyclicGroup 4) = 2^2
      1 (one_le_min_two_factorization_two h)

  The "exactly the same" part doesn't mean the proof terms need to be literally identical text. By proof irrelevance, any two proofs of Nat.card (CyclicGroup q) = q^1 are definitionally equal, so it doesn't
  matter if one comes from by rw [card_cyclicGroup, pow_one] and the other from by simp [card_cyclicGroup]. They're the same.

  The reason this is fast and the earlier attempt wasn't: when Lean sees sdpCanonicalAction ... (CyclicGroup q) (CyclicGroup 4) ... and canonicalC4OnCqAction ... (which unfolds to the exact same
  sdpCanonicalAction ... (CyclicGroup q) (CyclicGroup 4) ...), it only needs to check argument equality by congruence — same head function, same N, same K, all proof args equal by proof irrelevance. It never
  unfolds sdpCanonicalAction's body (which is where Classical.choice and the expensive stuff lives).

  The earlier timeout happened because I was comparing sdpCanonicalAction applied to (↥Q, ↥K) with sdpCanonicalAction applied to (CyclicGroup q, CyclicGroup 4) — different N and K type arguments — so Lean
  couldn't use congruence and was forced to unfold the body to check equality."

Claude gave up on me mid refactor so I want you to finish the job.

Essentially, the thing is, if you check git history, you would see use of horrible horrible bridges because the semi direct product between a cyclic p and q group classification result was stated in a bad way, which gave rose to the need for a lot of adaptors/bridges. Now its been improved and so it should be quite nice. Could you please make the code for 4q compile and make sure its neat and clean?

Here's the math (rough notes):

Lemma 1: For groups of order $p^2 q$, there must exist either a normal Sylow $p$-group or a normal Sylow $q$-group.

Lemma 2: If $K$ is a cyclic $p$-group, two homomorphisms $f, g : K → \text{Aut}(H)$ define
isomorphic semi-direct products $H \rtimes K$ if they have the same image.

Lemma 3: Semi-direct products $C_{q^n} ⋊ C_{p^m}$ ($q$ odd prime, $p ≠ q$) are classified up to
isomorphism by $r ∈ \{0, …, \min(m, d)\}$ where $d = v_p(q - 1)$, giving $\min(m, d) + 1$ classes. Every action $f$ belongs to exactly one class.

Lemma 4: If $|K| = p$, then two homomorphisms $f, g : K → \text{Aut}(H)$ define isomorphic semi-direct products $H \rtimes K$ if their images are conjugate subgroups of $\text{Aut}(H)$.

# Case 1: There exists a normal Sylow $p$-group

$|N| = p^2$ and by Schur-zassenhaus theorem there exists a complementary subgroup $U$ of $G$ with $|U| = q$ such that $G \cong N \rtimes U$. We have $U \cong C_q$.

## Case 1.1: $N \cong C_{p^2}$
$\text{Aut}(N) = C_{p(p-1)}$
Let $\phi : C_q \rightarrow C_{p(p-1)}$. Done by Lemma 3 when $p$ is odd. When $p = 2$, we get $\phi : C_q \rightarrow C_2$ so only trivial homomorphism is possible as $q \nmid 2$ for odd $q$.

(For $p = 2$, only trivial homomorphism)
(For $q = 2$, then we get 2 homomorphisms)

$C_{p^2 q}$ always
$C_{p^2} \rtimes C_q$ if $q \mid p - 1$.
## Case 1.2: $N \cong C_p \times C_p$
$\text{Aut}(N) \cong \text{GL}_2(p)$
Let $\phi : C_q \rightarrow \text{GL}_2(p)$

???
### Special case $p = 2$
$GL_2(2) \cong S_3$
Either $\phi$ is a trivial action, or its image is a subgroup of order $q$ in $S_3$, hence $q = 3$. Only one subgroup of order $3$ so gives one group.

### Special case $q = 2$
$\text{Im}(\phi)$ is trivial or $\langle A \rangle$ where $A \in GL_2(p)$ is of order $2$. By lemma 4 we only need to look at similarity classes and pick a representative from each class. Then $A^2 = 1$ implies minimal polynomial of $A$ divides $x^2 - 1$ and since $p$ is odd, it has distinct factors hence diagonalisable. Hence each such matrix similar to a diagonal matrix $\text{diag}(\alpha, \beta)$ such that $\alpha^2 = 1, \beta^2 = 1$ hence $\alpha = \pm 1, \beta = \pm 1$ ($t^2 - 1$ has at most $2$ roots because we are over a field). So $\alpha = \beta = 1$ is order $1$ so we don't consider it. We can check there are two similarity classes $\alpha = 1, \beta = -1$ (or switch them) and one for $\alpha = \beta = -1$.

Side note: Lemma 4 not needed because the general condition $\phi_2(u) = \alpha \circ \phi_1(\beta(u)) \circ \alpha^{-1}$ becomes equivalent to similarity because $\text{Aut}(C_q) = C_1$ as $q = 2$ so $\beta = \text{id}$ and that $\phi_1, \phi_2$ are determined by where they send the generator as $U$ is cyclic.

C_3 -> Aut(G)

# Case 2: There exists a normal Sylow $q$-group

$|N| = q$ and by Schur-zassenhaus theorem there exists a complementary subgroup $U$ of $G$ with $|U| = p^2$ such that $G \cong N \rtimes U$. We have $N \cong C_q$. $\text{Aut}(N) \cong C_{q-1}$.

## Case 2.1 $U \cong C_{p^2}$
Let $\phi : C_{p^2} \rightarrow C_{q-1}$. Done by Lemma 3 when $q$ is odd. Only trivial homomorphism possible when $q = 2$.

(When $p = 2$, have to see if $q$ is $1$ mod $4$)
(When $q = 2$, we get only trivial homomorphism)

$C_{p^2q}$ always
$C_q \rtimes_1 C_{p^2}$ if $v_p(q - 1) \geq 1$
$C_q \rtimes_2 C_{p^2}$ if $v_p(q - 1) \geq 2$
## Case 2.2 $U \cong C_p \times C_p$
Let $\phi : C_p \times C_p \rightarrow C_{q-1}$

$\text{Im}(\phi) \cong (C_p \times C_p) / \ker(\phi)$ so order of image is either $1, p$ or $p^2$. If it is $1$, then we have the trivial homomorphism giving $G \cong C_p \times C_p \times C_{q}$. If it is $p^2$, then we must have $\ker(\phi) = 1$ but $\text{Im}(\phi)$ would then be a non-cyclic subgroup of the cyclic group $C_{q-1}$, impossible. If it is $p$, then ?????

(When $p = 2$, one trivial and one non-trivial?)
(When $q = 2$, we get only trivial homomorphism)

$C_p \times C_p \times C_q$
$C_q \rtimes (C_p \times C_p)$ if $p \mid q - 1$

%%FILE numina/blueprints/4q-group-classification-for-q-3-and-q-is-3-mod-4/4q-group-classification-for-q-3-and-q-is-3-mod-4.tex
% Blueprint: 4q group classification for q > 3 and q is 3 mod 4
% Created by Numina

\begin{document}

\section{Setup and assumptions}

We classify (partially) finite groups $G$ of order $4q$ where $q$ is a prime with
$q > 3$ and $q \equiv 3 \pmod 4$. We work with the additive shorthand $p = 2$ so
that $|G| = p^2 q = 4q$. The two isomorphism types we will exhibit are the
cyclic group $C_{4q}$ and the elementary direct product
$C_2 \times C_2 \times C_q$.

\begin{remark}
    The complete classification of groups of order $4q$ for $q \equiv 3
    \pmod 4$ contains four isomorphism types, the two above together with
    $C_q \rtimes C_4$ (nontrivial) and $C_q \rtimes (C_2 \times C_2) \cong
    D_{2q}$. The current target theorem in this blueprint asserts only the
    disjunction $G \cong C_{4q}$ or $G \cong C_2 \times C_2 \times C_q$; the
    nontrivial semidirect-product cases (Case 2 of the proof) are documented
    here but treated as a residual TODO.
\end{remark}

\section{Sylow structure}

\begin{lemma}[A group of order $p^2 q$ has a normal Sylow subgroup]
    \label{lem:p2q-normal-sylow}
    Let $p, q$ be distinct primes and let $G$ be a finite group of order
    $p^2 q$. Then either the number of Sylow $p$-subgroups of $G$ is $1$ (so
    the unique Sylow $p$-subgroup is normal) or the number of Sylow
    $q$-subgroups of $G$ is $1$ (so the unique Sylow $q$-subgroup is normal).
\end{lemma}
\begin{proof}
    Standard Sylow counting. Already established in the repository as
    \texttt{p2q\_group\_has\_normal\_sylow\_subgroup}.
\end{proof}

\begin{lemma}[Sylow cardinality and index]
    \label{lem:sylow-card-index}
    Let $p, q$ be distinct primes, $a, b \in \mathbb{N}$, and let $G$ be a
    finite group with $|G| = p^a q^b$. Then every Sylow $p$-subgroup of $G$
    has order $p^a$ and index $q^b$.
\end{lemma}
\begin{proof}
    From $|G| = p^a q^b$ and coprimality of $p^a$ and $q^b$, the multiplicity
    of $p$ in $|G|$ is $a$, so every Sylow $p$-subgroup has order $p^a$ by
    Mathlib's \texttt{Sylow.card\_eq\_multiplicity}. The index is then
    $|G|/p^a = q^b$ by Lagrange. Already established as
    \texttt{sylow\_card\_eq} and \texttt{sylow\_index\_eq}.
\end{proof}

\section{Background results on small groups}

\begin{lemma}[Cyclic classification of prime order]
    \label{lem:prime-classification}
    If $|G| = n$ with $n$ prime, then $G \cong C_n$.
\end{lemma}
\begin{proof}
    A group of prime order has a generator (any non-identity element), hence
    is cyclic. Two cyclic groups of equal finite order are isomorphic via
    Mathlib's \texttt{mulEquivOfCyclicCardEq}. Already established as
    \texttt{prime\_classification}.
\end{proof}

\begin{lemma}[Classification of groups of order $p^2$]
    \label{lem:p-squared-classification}
    For any prime $p$, a group of order $p^2$ is isomorphic to either
    $C_{p^2}$ or $C_p \times C_p$.
\end{lemma}
\begin{proof}
    Any group of order $p^2$ is abelian; either it is cyclic (and isomorphic
    to $C_{p^2}$) or it has exponent $p$, in which case it is a $2$-dimensional
    $\mathbb{F}_p$-vector space. Already established as
    \texttt{p\_squared\_classification}.
\end{proof}

\begin{lemma}[Aut$(C_{p^2}) \cong C_{p(p-1)}$]
    \label{lem:aut-cyclic-p2}
    For any prime $p$, $\mathrm{MulAut}(C_{p^2}) \cong C_{p(p-1)}$.
\end{lemma}
\begin{proof}
    $\mathrm{MulAut}(C_{p^2}) \cong (\mathbb{Z}/p^2 \mathbb{Z})^\times$,
    which is cyclic of order $\varphi(p^2) = p(p-1)$. Already established as
    \texttt{aut\_of\_cyclic\_p2}.
\end{proof}

\begin{lemma}[Aut$(C_p \times C_p) \cong \mathrm{GL}_2(\mathbb{F}_p)$]
    \label{lem:aut-CpCp}
    For any prime $p$, $\mathrm{MulAut}(C_p \times C_p) \cong
    \mathrm{GL}_2(\mathbb{F}_p)$.
\end{lemma}
\begin{proof}
    $C_p \times C_p$ is a $2$-dimensional $\mathbb{F}_p$-vector space and
    every additive automorphism is automatically $\mathbb{F}_p$-linear.
    Already established as \texttt{aut\_of\_CpCp}.
\end{proof}

\begin{lemma}[$\mathrm{GL}_2(\mathbb{F}_2) \cong S_3$]
    \label{lem:GL2F2-isoS3}
    $\mathrm{GL}_2(\mathbb{F}_2)$ is isomorphic to the dihedral group of
    order $6$ (equivalently, $S_3$).
\end{lemma}
\begin{proof}
    A computation: $|\mathrm{GL}_2(\mathbb{F}_2)| = 6$ and
    $\mathrm{GL}_2(\mathbb{F}_2)$ is non-abelian, so by the classification of
    groups of order $6$ it must be isomorphic to $D_3$. Already established
    as \texttt{GL2F2\_isoS3}.
\end{proof}

\section{Triviality criteria for actions}

\begin{lemma}[Trivial homomorphism from coprime cardinalities]
    \label{lem:eq-one-of-coprime-card}
    Let $G, H$ be finite groups with $\gcd(|G|, |H|) = 1$. Then every group
    homomorphism $\varphi : G \to H$ is the trivial homomorphism $\varphi = 1$.
\end{lemma}
\begin{proof}
    The image $\varphi(G)$ is a subgroup of $H$ whose order divides both
    $|G|$ and $|H|$, hence divides $\gcd(|G|,|H|) = 1$. Therefore
    $|\varphi(G)| = 1$, so $\varphi(G) = \{1\}$, so $\varphi = 1$.
    Already established as \texttt{eq\_one\_of\_coprime\_card}.
\end{proof}

\begin{lemma}[Semidirect product with trivial action is a direct product]
    \label{lem:sdp-trivial-action}
    Let $N, K$ be groups and let $\varphi : K \to \mathrm{Aut}(N)$ be a
    homomorphism with $\varphi = 1$. Then $N \rtimes_\varphi K \cong N \times K$.
\end{lemma}
\begin{proof}
    The underlying set of $N \rtimes_\varphi K$ is $N \times K$ and the
    multiplication
    $(n_1, k_1)(n_2, k_2) = (n_1 \cdot \varphi(k_1)(n_2),\ k_1 k_2)$
    reduces to $(n_1 n_2, k_1 k_2)$ when $\varphi \equiv 1$. Hence the
    identity pair map is a group isomorphism. Already established as
    \texttt{SemidirectProduct.mulEquivOfTrivialAction}.
\end{proof}

\begin{lemma}[Coprime cyclic direct product]
    \label{lem:cyclic-prod}
    For positive integers $m, n$ with $\gcd(m, n) = 1$,
    $C_m \times C_n \cong C_{mn}$.
\end{lemma}
\begin{proof}
    $C_m \times C_n$ is cyclic (the direct product of two cyclic groups of
    coprime orders is cyclic, since one can take the product of generators
    and apply the Chinese Remainder Theorem on orders), and has cardinality
    $mn$. Already established as \texttt{CyclicGroup.prodMulEquiv}.
\end{proof}

\section{Existence of complements}

\begin{lemma}[Schur--Zassenhaus complement]
    \label{lem:schur-zassenhaus}
    Let $G$ be a finite group and $N \trianglelefteq G$ a normal subgroup
    such that $\gcd(|N|, [G : N]) = 1$. Then there exists a subgroup
    $K \leq G$ with $N \cap K = \{1\}$ and $N K = G$ (a right complement to
    $N$ in $G$).
\end{lemma}
\begin{proof}
    The Schur--Zassenhaus theorem, available in Mathlib as
    \texttt{Subgroup.exists\_right\_complement'\_of\_coprime}.
\end{proof}

\begin{lemma}[Semidirect product from a complement]
    \label{lem:sdp-from-complement}
    Let $G$ be a group, $N \trianglelefteq G$ a normal subgroup, and
    $K \leq G$ a complement of $N$ (so $N \cap K = \{1\}$ and $NK = G$).
    Then $G \cong N \rtimes_\varphi K$, where
    $\varphi : K \to \mathrm{Aut}(N)$ is the conjugation action.
\end{lemma}
\begin{proof}
    Mathlib provides this isomorphism as
    \texttt{SemidirectProduct.mulEquivSubgroup}.
\end{proof}

\section{The main classification (target theorem)}

\begin{theorem}[Partial classification of groups of order $4q$]
    \label{thm:classification-4q}
    \uses{lem:p2q-normal-sylow, lem:sylow-card-index,
          lem:prime-classification, lem:p-squared-classification,
          lem:aut-cyclic-p2, lem:aut-CpCp, lem:GL2F2-isoS3,
          lem:eq-one-of-coprime-card, lem:sdp-trivial-action,
          lem:cyclic-prod, lem:schur-zassenhaus, lem:sdp-from-complement,
          lem:case1-c4, lem:case1-c2c2}
    Let $q$ be a prime with $q > 3$ and $q \equiv 3 \pmod 4$, and let $G$ be
    a group with $|G| = 4q$. Then
    \[
        G \cong C_{4q} \quad \text{or} \quad G \cong C_2 \times C_2 \times C_q.
    \]
    (This is a partial classification: when the unique Sylow $q$-subgroup is
    normal and the conjugation action is nontrivial, $G$ falls outside this
    disjunction; see the Case 2 discussion below for the residual cases.)
\end{theorem}
\begin{proof}
    \uses{lem:p2q-normal-sylow, lem:case1-c4, lem:case1-c2c2,
          lem:schur-zassenhaus, lem:sdp-from-complement,
          lem:sylow-card-index}
    By Lemma~\ref{lem:p2q-normal-sylow} applied with $p = 2$, either the
    Sylow $2$-subgroup or the Sylow $q$-subgroup of $G$ is normal.

    \emph{Case 1: Normal Sylow $2$-subgroup $P$ with $|P| = 4$ and $[G:P] = q$.}
    By Lemma~\ref{lem:sylow-card-index} we have $|P| = 4$ and $[G:P] = q$. By
    Lemma~\ref{lem:schur-zassenhaus} (with $\gcd(4, q) = 1$ since $q$ is an
    odd prime) there is a complement $K \leq G$ with $|K| = q$. By
    Lemma~\ref{lem:sdp-from-complement}, $G \cong P \rtimes_\varphi K$ for
    some $\varphi : K \to \mathrm{Aut}(P)$. By
    Lemma~\ref{lem:prime-classification}, $K \cong C_q$. By
    Lemma~\ref{lem:p-squared-classification}, $P \cong C_4$ or
    $P \cong C_2 \times C_2$. Apply Lemma~\ref{lem:case1-c4} or
    Lemma~\ref{lem:case1-c2c2} accordingly to conclude
    $G \cong C_{4q}$ or $G \cong C_2 \times C_2 \times C_q$.

    \emph{Case 2: Normal Sylow $q$-subgroup $N$ with $|N| = q$ and $[G:N] = 4$.}
    By Lemma~\ref{lem:sylow-card-index}, $|N| = q$ and $[G:N] = 4$. By
    Lemma~\ref{lem:schur-zassenhaus} there is a complement $K$ with
    $|K| = 4$ and by Lemma~\ref{lem:sdp-from-complement},
    $G \cong N \rtimes_\varphi K$ with $N \cong C_q$ and $K \cong C_4$ or
    $C_2 \times C_2$. When $\varphi$ is trivial, $G \cong C_q \times K$, which
    gives $G \cong C_{4q}$ (if $K \cong C_4$, by
    Lemma~\ref{lem:cyclic-prod}) or $G \cong C_2 \times C_2 \times C_q$ (if
    $K \cong C_2 \times C_2$). When $\varphi$ is nontrivial, the resulting
    semidirect product is not isomorphic to either group in the disjunction;
    these subcases lie outside the current target theorem and are listed as
    TODO in the discussion of Case 2 below.
\end{proof}

\section{Case 1: normal Sylow $2$-subgroup}

The two subcases of Case 1 depend on whether the Sylow $2$-subgroup $P$ (of
order $4$) is isomorphic to $C_4$ or to $C_2 \times C_2$.

\begin{lemma}[Case 1.1: $P \cong C_4$ gives $G \cong C_{4q}$]
    \label{lem:case1-c4}
    \uses{lem:aut-cyclic-p2, lem:eq-one-of-coprime-card,
          lem:sdp-trivial-action, lem:cyclic-prod}
    Let $q$ be an odd prime, let $P$ be a group with $P \cong C_4$, let
    $K$ be a group with $K \cong C_q$, and let
    $\varphi : K \to \mathrm{MulAut}(P)$ be any homomorphism. Then
    $P \rtimes_\varphi K \cong C_{4q}$.
\end{lemma}
\begin{proof}
    \uses{lem:aut-cyclic-p2, lem:eq-one-of-coprime-card,
          lem:sdp-trivial-action, lem:cyclic-prod}
    By Lemma~\ref{lem:aut-cyclic-p2} applied to $p = 2$ and using
    $P \cong C_4 = C_{2^2}$, we have
    $\mathrm{MulAut}(P) \cong C_{2 \cdot 1} = C_2$, so
    $|\mathrm{MulAut}(P)| = 2$. Since $q > 2$ is an odd prime, $\gcd(q, 2) = 1$.
    By Lemma~\ref{lem:eq-one-of-coprime-card} applied to
    $\varphi : K \to \mathrm{MulAut}(P)$, $\varphi$ is trivial. By
    Lemma~\ref{lem:sdp-trivial-action},
    $P \rtimes_\varphi K \cong P \times K \cong C_4 \times C_q$. Finally, by
    Lemma~\ref{lem:cyclic-prod} with $\gcd(4, q) = 1$, we conclude
    $C_4 \times C_q \cong C_{4q}$.
\end{proof}

\begin{lemma}[Case 1.2: $P \cong C_2 \times C_2$ gives $G \cong C_2 \times C_2 \times C_q$]
    \label{lem:case1-c2c2}
    \uses{lem:aut-CpCp, lem:GL2F2-isoS3, lem:eq-one-of-coprime-card,
          lem:sdp-trivial-action}
    Let $q$ be a prime with $q > 3$, let $P$ be a group with
    $P \cong C_2 \times C_2$, let $K$ be a group with $K \cong C_q$, and let
    $\varphi : K \to \mathrm{MulAut}(P)$ be any homomorphism. Then
    $P \rtimes_\varphi K \cong C_2 \times C_2 \times C_q$.
\end{lemma}
\begin{proof}
    \uses{lem:aut-CpCp, lem:GL2F2-isoS3, lem:eq-one-of-coprime-card,
          lem:sdp-trivial-action}
    By Lemma~\ref{lem:aut-CpCp}, $\mathrm{MulAut}(C_2 \times C_2) \cong
    \mathrm{GL}_2(\mathbb{F}_2)$, and by Lemma~\ref{lem:GL2F2-isoS3},
    $\mathrm{GL}_2(\mathbb{F}_2) \cong D_3$ (a group of order $6$). Hence
    $|\mathrm{MulAut}(P)| = 6 = 2 \cdot 3$. Since $q$ is prime with $q > 3$,
    $q$ is coprime to both $2$ and $3$, hence $\gcd(q, 6) = 1$. By
    Lemma~\ref{lem:eq-one-of-coprime-card},
    $\varphi : K \to \mathrm{MulAut}(P)$ is trivial. By
    Lemma~\ref{lem:sdp-trivial-action},
    $P \rtimes_\varphi K \cong P \times K \cong (C_2 \times C_2) \times C_q
    \cong C_2 \times C_2 \times C_q$ (the last step is reassociation of the
    direct product, an instance of \texttt{MulEquiv.prodAssoc}).
\end{proof}

\section{Case 2 (residual): normal Sylow $q$-subgroup}

In Case 2 the unique Sylow $q$-subgroup $N$ is normal with $|N| = q$, and a
Schur--Zassenhaus complement $K$ has $|K| = 4$. The conjugation action gives
$\varphi : K \to \mathrm{MulAut}(N)$. Since $N \cong C_q$,
$\mathrm{MulAut}(N) \cong C_{q-1}$.

For $q \equiv 3 \pmod 4$, $v_2(q-1) = 1$, so a homomorphism from any group of
order $4$ into $C_{q-1}$ has image of order $1$ or $2$. The image-of-order-$1$
case is the trivial action and reduces to the direct product
$N \times K$, contributing
\begin{itemize}
    \item $C_q \times C_4 \cong C_{4q}$ when $K \cong C_4$ (already covered by
        the Case 1 conclusion via Lemma~\ref{lem:cyclic-prod}), and
    \item $C_q \times (C_2 \times C_2) \cong C_2 \times C_2 \times C_q$ when
        $K \cong C_2 \times C_2$.
\end{itemize}

The image-of-order-$2$ (nontrivial action) cases give the additional
isomorphism types
\[
    C_q \rtimes C_4 \quad \text{and} \quad C_q \rtimes (C_2 \times C_2) \cong D_{2q},
\]
which fall outside the disjunction $G \cong C_{4q} \lor G \cong C_2 \times
C_2 \times C_q$ in the target theorem. These subcases are therefore left as
TODO; the current target theorem statement is not provable in full as stated
without strengthening either the hypotheses (excluding Case 2 nontrivial
actions) or the conclusion (adding the two further isomorphism classes).

\end{document}


%%FILE numina/blueprints/4q-group-classification-for-q-3-and-q-is-3-mod-4/source/4q-group-classification-for-q-3-and-q-is-3-mod-4-source.tex
Lemma 1: For groups of order $p^2 q$, there must exist either a normal Sylow $p$-group or a normal Sylow $q$-group.

Lemma 2: If $K$ is a cyclic $p$-group, two homomorphisms $f, g : K → \text{Aut}(H)$ define
isomorphic semi-direct products $H \rtimes K$ if they have the same image.

Lemma 3: Semi-direct products $C_{q^n} ⋊ C_{p^m}$ ($q$ odd prime, $p ≠ q$) are classified up to
isomorphism by $r ∈ \{0, …, \min(m, d)\}$ where $d = v_p(q - 1)$, giving $\min(m, d) + 1$ classes. Every action $f$ belongs to exactly one class.

Lemma 4: If $|K| = p$, then two homomorphisms $f, g : K → \text{Aut}(H)$ define isomorphic semi-direct products $H \rtimes K$ if their images are conjugate subgroups of $\text{Aut}(H)$.

$|N| = p^2$ and by Schur-zassenhaus theorem there exists a complementary subgroup $U$ of $G$ with $|U| = q$ such that $G \cong N \rtimes U$. We have $U \cong C_q$.

## Case 1.1: $N \cong C_{p^2}$
$\text{Aut}(N) = C_{p(p-1)}$
Let $\phi : C_q \rightarrow C_{p(p-1)}$. Done by Lemma 3 when $p$ is odd. When $p = 2$, we get $\phi : C_q \rightarrow C_2$ so only trivial homomorphism is possible as $q \nmid 2$ for odd $q$.

(For $p = 2$, only trivial homomorphism)
(For $q = 2$, then we get 2 homomorphisms)

## Case 1.2: $N \cong C_p \times C_p$
$\text{Aut}(N) \cong \text{GL}_2(p)$
Let $\phi : C_q \rightarrow \text{GL}_2(p)$
### Special case $p = 2$
$GL_2(2) \cong S_3$
Either $\phi$ is a trivial action, or its image is a subgroup of order $q$ in $S_3$, hence $q = 3$. Only one subgroup of order $3$ so gives one group.

### Special case $q = 2$
$\text{Im}(\phi)$ is trivial or $\langle A \rangle$ where $A \in GL_2(p)$ is of order $2$. By lemma 4 we only need to look at similarity classes and pick a representative from each class. Then $A^2 = 1$ implies minimal polynomial of $A$ divides $x^2 - 1$ and since $p$ is odd, it has distinct factors hence diagonalisable. Hence each such matrix similar to a diagonal matrix $\text{diag}(\alpha, \beta)$ such that $\alpha^2 = 1, \beta^2 = 1$ hence $\alpha = \pm 1, \beta = \pm 1$ ($t^2 - 1$ has at most $2$ roots because we are over a field). So $\alpha = \beta = 1$ is order $1$ so we don't consider it. We can check there are two similarity classes $\alpha = 1, \beta = -1$ (or switch them) and one for $\alpha = \beta = -1$.

Side note: Lemma 4 not needed because the general condition $\phi_2(u) = \alpha \circ \phi_1(\beta(u)) \circ \alpha^{-1}$ becomes equivalent to similarity because $\text{Aut}(C_q) = C_1$ as $q = 2$ so $\beta = \text{id}$ and that $\phi_1, \phi_2$ are determined by where they send the generator as $U$ is cyclic.

# Case 2: There exists a normal Sylow $q$-group

$|N| = q$ and by Schur-zassenhaus theorem there exists a complementary subgroup $U$ of $G$ with $|U| = p^2$ such that $G \cong N \rtimes U$. We have $N \cong C_q$. $\text{Aut}(N) \cong C_{q-1}$.

## Case 2.1 $U \cong C_{p^2}$
Let $\phi : C_{p^2} \rightarrow C_{q-1}$. Done by Lemma 3 when $q$ is odd. Only trivial homomorphism possible when $q = 2$.

(When $p = 2$, have to see if $q$ is $1$ mod $4$)
(When $q = 2$, we get only trivial homomorphism)

## Case 2.2 $U \cong C_p \times C_p$
Let $\phi : C_p \times C_p \rightarrow C_{q-1}$

$\text{Im}(\phi) \cong (C_p \times C_p) / \ker(\phi)$ so order of image is either $1, p$ or $p^2$. If it is $1$, then we have the trivial homomorphism giving $G \cong C_p \times C_p \times C_{q}$. If it is $p^2$, then we must have $\ker(\phi) = 1$ but $\text{Im}(\phi)$ would then be a non-cyclic subgroup of the cyclic group $C_{q-1}$, impossible. If it is $p$, then ?????

(When $p = 2$, one trivial and one non-trivial?)
(When $q = 2$, we get only trivial homomorphism)
# p = 2, q > 3, q is 1 mod 4
1 (1.1) + 1 (1.2) + 3 (2.1) + 2 (2.2) = 6

According to formula: 5, we get $C_{4q}$ again from case 2.1 (trivial homomorphism)
$C_{4q}$ (1.1)
$C_2 \times C_2 \times C_q$ (1.2)
$C_q \rtimes_1 C_4$ (2.1)
$C_q \rtimes_2 C_4$ (2.1)
$C_q \rtimes (C_2 \times C_2)$ (2.2)
# p = 2, q > 3, q is 3 mod 4
1 (1.1) + 1 (1.2) + 2 (2.1) + 2 (2.2) = 6

$C_{4q}$ (1.1)
$C_2 \times C_2 \times C_q$ (1.2)
$C_q \rtimes C_4$ (2.1)
$C_q \rtimes (C_2 \times C_2)$ (2.2)

According to formula: 4, we get $C_{p^2 q}$ again from case 2.1 (trivial homomorphism)
# p = 2, q = 3
1 (1.1) + 2 (1.2) + 2 (2.1) + 2 (2.2) = 7
There are only 5 so some 2 are isomorphic here

$C_{12}$ (1.1)
$C_2 \times C_2 \times C_3$ (1.2)
$(C_2 \times C_2) \rtimes C_3$ (1.2)
$C_3 \rtimes C_4$ (2.1)
$C_3 \rtimes (C_2 \times C_2)$ (2.2)
# q = 2
2 (1.1), 2 (1.2), 1 (2.1, $C_{2p^2}$ repeated), 1 (2.2) total = 5

%%FILE numina/blueprints/classification-of-groups-of-order-2-p2/classification-of-groups-of-order-2-p2.tex
% Blueprint: Classification of groups of order 2 p^2
% Created by Numina

\begin{document}

% Your blueprint will go here.

\end{document}


%%FILE numina/blueprints/classification-of-groups-of-order-2-p2/source/classification-of-groups-of-order-2-p2-source.tex
Lemma 1: For groups of order $p^2 q$, there must exist either a normal Sylow $p$-group or a normal Sylow $q$-group.

Lemma 2: If $K$ is a cyclic $p$-group, two homomorphisms $f, g : K → \text{Aut}(H)$ define
isomorphic semi-direct products $H \rtimes K$ if they have the same image.

Lemma 3: Semi-direct products $C_{q^n} ⋊ C_{p^m}$ ($q$ odd prime, $p ≠ q$) are classified up to
isomorphism by $r ∈ \{0, …, \min(m, d)\}$ where $d = v_p(q - 1)$, giving $\min(m, d) + 1$ classes. Every action $f$ belongs to exactly one class.

Lemma 4: If $|K| = p$, then two homomorphisms $f, g : K → \text{Aut}(H)$ define isomorphic semi-direct products $H \rtimes K$ if their images are conjugate subgroups of $\text{Aut}(H)$.

# Case 1: There exists a normal Sylow $p$-group

$|N| = p^2$ and by Schur-zassenhaus theorem there exists a complementary subgroup $U$ of $G$ with $|U| = q$ such that $G \cong N \rtimes U$. We have $U \cong C_q$.

## Case 1.1: $N \cong C_{p^2}$
$\text{Aut}(N) = C_{p(p-1)}$
Let $\phi : C_q \rightarrow C_{p(p-1)}$. Done by Lemma 3 when $p$ is odd. When $p = 2$, we get $\phi : C_q \rightarrow C_2$ so only trivial homomorphism is possible as $q \nmid 2$ for odd $q$.

(For $p = 2$, only trivial homomorphism)
(For $q = 2$, then we get 2 homomorphisms)

$C_{p^2 q}$ always
$C_{p^2} \rtimes C_q$ if $q \mid p - 1$.
## Case 1.2: $N \cong C_p \times C_p$
$\text{Aut}(N) \cong \text{GL}_2(p)$
Let $\phi : C_q \rightarrow \text{GL}_2(p)$

???
### Special case $p = 2$
$GL_2(2) \cong S_3$
Either $\phi$ is a trivial action, or its image is a subgroup of order $q$ in $S_3$, hence $q = 3$. Only one subgroup of order $3$ so gives one group.

### Special case $q = 2$
$\text{Im}(\phi)$ is trivial or $\langle A \rangle$ where $A \in GL_2(p)$ is of order $2$. By lemma 4 we only need to look at similarity classes and pick a representative from each class. Then $A^2 = 1$ implies minimal polynomial of $A$ divides $x^2 - 1$ and since $p$ is odd, it has distinct factors hence diagonalisable. Hence each such matrix similar to a diagonal matrix $\text{diag}(\alpha, \beta)$ such that $\alpha^2 = 1, \beta^2 = 1$ hence $\alpha = \pm 1, \beta = \pm 1$ ($t^2 - 1$ has at most $2$ roots because we are over a field). So $\alpha = \beta = 1$ is order $1$ so we don't consider it. We can check there are two similarity classes $\alpha = 1, \beta = -1$ (or switch them) and one for $\alpha = \beta = -1$.

Side note: Lemma 4 not needed because the general condition $\phi_2(u) = \alpha \circ \phi_1(\beta(u)) \circ \alpha^{-1}$ becomes equivalent to similarity because $\text{Aut}(C_q) = C_1$ as $q = 2$ so $\beta = \text{id}$ and that $\phi_1, \phi_2$ are determined by where they send the generator as $U$ is cyclic.

# Case 2: There exists a normal Sylow $q$-group

$|N| = q$ and by Schur-zassenhaus theorem there exists a complementary subgroup $U$ of $G$ with $|U| = p^2$ such that $G \cong N \rtimes U$. We have $N \cong C_q$. $\text{Aut}(N) \cong C_{q-1}$.

## Case 2.1 $U \cong C_{p^2}$
Let $\phi : C_{p^2} \rightarrow C_{q-1}$. Done by Lemma 3 when $q$ is odd. Only trivial homomorphism possible when $q = 2$.

(When $p = 2$, have to see if $q$ is $1$ mod $4$)
(When $q = 2$, we get only trivial homomorphism)

$C_{p^2q}$ always
$C_q \rtimes_1 C_{p^2}$ if $v_p(q - 1) \geq 1$
$C_q \rtimes_2 C_{p^2}$ if $v_p(q - 1) \geq 2$
## Case 2.2 $U \cong C_p \times C_p$
Let $\phi : C_p \times C_p \rightarrow C_{q-1}$

$\text{Im}(\phi) \cong (C_p \times C_p) / \ker(\phi)$ so order of image is either $1, p$ or $p^2$. If it is $1$, then we have the trivial homomorphism giving $G \cong C_p \times C_p \times C_{q}$. If it is $p^2$, then we must have $\ker(\phi) = 1$ but $\text{Im}(\phi)$ would then be a non-cyclic subgroup of the cyclic group $C_{q-1}$, impossible. If it is $p$, then ?????

(When $p = 2$, one trivial and one non-trivial?)
(When $q = 2$, we get only trivial homomorphism)

$C_p \times C_p \times C_q$
$C_q \rtimes (C_p \times C_p)$ if $p \mid q - 1$

---

\begin{table}[htbp]
\centering
\caption{Classification of Groups of Order $2p^2$ ($q = 2$, $p$ is odd)}
\label{tab:groups_p2q_classification}
\begin{tabular}{ll p{3.2cm} p{6cm}}
\toprule
\textbf{Subcase} & \textbf{Group Structure ($G$)} & \textbf{Condition for existence} & \textbf{Proof Details / References} \\
\midrule
\quad $N \cong C_{p^2}$ & $C_{2p^2}$ & Always & Trivial action ($\phi$ is trivial). \\
\quad $N \cong C_{p^2}$ & $C_{p^2} \rtimes C_2$ & Always & Non-trivial action; see \ref{Classification of semi-direct products between Cyclic p-group and q-groups}. \\
\addlinespace
\quad $N \cong C_p \times C_p$ & $C_p \times C_p \times C_2$ & Always & Trivial action ($\phi$ is trivial) \\
\quad $N \cong C_p \times C_p$ & $(C_p \times C_p) \rtimes_1 C_2$ & Always & Image of action is conjugate to $\text{diag}(1, -1)$ over $GL_2(\mathbb{F}_p)$ \\
\quad $N \cong C_p \times C_p$ & $(C_p \times C_p) \rtimes_2 C_2$ & Always & Image of action is conjugate to $\text{diag}(-1, -1)$ over $GL_2(\mathbb{F}_p)$ \\
\midrule
\quad $U \cong C_{p^2}$ & $C_{2p^2}$ & Always & Trivial action ($\phi$ is trivial). \\
\addlinespace
\quad $U \cong C_p \times C_p$ & $C_p \times C_p \times C_2$ & Always & Trivial action ($\phi$ is trivial). \\
\bottomrule
\end{tabular}
\end{table}

%%FILE numina/blueprints/conjugacy-classes-of-elements-of-order-two-of-gl-2f-p/conjugacy-classes-of-elements-of-order-two-of-gl-2f-p.tex
% Blueprint: Conjugacy classes of elements of order two of GL_2(F_p)
% Created by Numina

\begin{document}

\section{Conjugacy classes of order-two elements of \texorpdfstring{$\mathrm{GL}_2(\mathbb{F}_p)$}{GL2(Fp)}}

Throughout this section, $p$ is an odd prime and $\mathbb{F}_p = \mathbb{Z}/p\mathbb{Z}$.

\begin{definition}
  \label{def:gl2-diag-one-neg-one}
  \lean{gl2Diag1NegOne}
  \leanfile{M2rGroup7/Lemmas/LinearAlgebraUtils.lean}
  \leanok
  The element $\operatorname{diag}(1, -1) \in \mathrm{GL}_2(\mathbb{F}_p)$.
\end{definition}

\begin{definition}
  \label{def:gl2-diag-neg-one-neg-one}
  \lean{gl2DiagNeg1Neg1}
  \leanfile{M2rGroup7/Lemmas/LinearAlgebraUtils.lean}
  \leanok
  The element $\operatorname{diag}(-1, -1) = -I \in \mathrm{GL}_2(\mathbb{F}_p)$.
\end{definition}

\begin{theorem}
  \label{thm:gl2-order-two-classification}
  \lean{gl2_order_two_classification}
  \leanfile{M2rGroup7/Lemmas/LinearAlgebraUtils.lean}
  \uses{def:gl2-diag-one-neg-one, def:gl2-diag-neg-one-neg-one}
  \leanok
  Let $p$ be an odd prime and let $A \in \mathrm{GL}_2(\mathbb{F}_p)$ be an element of order $2$.
  Then $A$ is conjugate in $\mathrm{GL}_2(\mathbb{F}_p)$ to either $\operatorname{diag}(1, -1)$ or to
  $\operatorname{diag}(-1, -1) = -I$.
\end{theorem}

\begin{proof}
  Let $M = A$ regarded as a matrix. Since $A$ has order $2$, we have $M^2 = I$ and $M \neq I$.

  If $M = -I$, then $A = \operatorname{diag}(-1, -1)$ is conjugate to itself by the identity, giving the second alternative.

  Otherwise $M \neq I$ and $M \neq -I$. Since $M^2 - I = (M - I)(M + I) = 0$ and $M \neq -I$, the matrix
  $M + I$ is nonzero, so it has a nonzero column $v_+$. From $(M - I)(M + I) = 0$ we get
  $M v_+ = v_+$, i.e.\ $v_+$ is a $1$-eigenvector. Similarly, since $M \neq I$, the matrix $M - I$ is
  nonzero, so it has a nonzero column $v_-$, and $(M + I)(M - I) = 0$ gives $M v_- = -v_-$, so $v_-$
  is a $(-1)$-eigenvector.

  Because $1 \neq -1$ in $\mathbb{F}_p$ (this is where we use $p \neq 2$), eigenvectors for distinct
  eigenvalues are linearly independent, so $\{v_+, v_-\}$ is a basis of $\mathbb{F}_p^2$. The change-of-basis
  matrix $P = [\,v_+ \mid v_-\,]$ lies in $\mathrm{GL}_2(\mathbb{F}_p)$ and satisfies
  $M P = P \cdot \operatorname{diag}(1, -1)$, i.e.\ $P^{-1} M P = \operatorname{diag}(1, -1)$. Hence
  $A$ is conjugate to $\operatorname{diag}(1, -1)$, giving the first alternative.
\end{proof}

\end{document}


%%FILE numina/blueprints/conjugacy-classes-of-elements-of-order-two-of-gl-2f-p/source/conjugacy-classes-of-elements-of-order-two-of-gl-2f-p-source.tex
(Part of a proof of the following theorem in LinearAlgebraUtils.lean)

theorem gl2_order_two_classification (hp2 : p ≠ 2) (A : GL (Fin 2) (ZMod p))
    (hA : orderOf A = 2) :
    IsConj A (gl2Diag1NegOne hp2) ∨ IsConj A (gl2DiagNeg1Neg1 : GL (Fin 2) (ZMod p))

...

-- Strategy: since M² = I, M ≠ I, M ≠ -I, the minimal polynomial of M divides
    -- (X-1)(X+1), which is squarefree (p odd). Both eigenvalues
    -- ±1 lie in 𝔽_p, so M is diagonalizable over 𝔽_p with eigenvalues {1, -1}.
    -- Concretely: since M ≠ -I, the matrix M + I ≠ 0, so pick a nonzero column v₊
    -- of M + I; it satisfies M * v₊ = v₊ (a 1-eigenvector, since (M-I)(M+I) = M²-I = 0).
    -- Similarly, since M ≠ I, pick a nonzero column v₋ of M - I satisfying M * v₋ = -v₋.
    -- Since 1 ≠ -1 (p ≠ 2), v₊ and v₋ are linearly independent.
    -- The matrix P = [v₊ | v₋] ∈ GL₂(𝔽_p) conjugates A to diag(1,-1).

%%FILE numina/blueprints/factoring-out-lemma-about-center-of-semi-direct-product/factoring-out-lemma-about-center-of-semi-direct-product.tex
% Blueprint: Factoring out lemma about center of semi direct product
% Created by Numina

\begin{document}

\section{Centre of a semidirect product of abelian groups}

The classification of semidirect products $C_{q^n}\rtimes C_{p^m}$ in
\texttt{CycPGroupClassification.lean} uses a description of the centre that, in
the existing proof, is stated for cyclic $p$- and $q$-groups but in fact only
depends on the abelianness of the two factors. We factor that lemma out into a
standalone statement for arbitrary abelian groups.

Throughout this section, $N$ and $K$ are (multiplicative) abelian groups, and
$f : K \to \operatorname{Aut}(N)$ is a group homomorphism. Write $N\rtimes_f K$
for the semidirect product with the operation
\[
  (x_1, y_1)(x_2, y_2) \;=\; \bigl(x_1\cdot f(y_1)(x_2),\; y_1 y_2\bigr).
\]

\begin{definition}\label{def:fixedPointsSubgroup}
\lean{fixedPointsSubgroup}
\leanfile{M2rGroup7/Lemmas/GroupTheoryLemmas.lean}
\leanok
The \emph{fixed-point subgroup} $\operatorname{Fix}(\operatorname{Im} f)\le N$
of $f$ is
\[
  \operatorname{Fix}(\operatorname{Im} f)
  \;=\; \{\, n \in N \;:\; f(h)(n) = n \text{ for every } h \in K \,\}.
\]
It is a subgroup because $f(h)$ is a group automorphism of $N$ for every
$h \in K$.
\end{definition}

\begin{lemma}\label{lem:mem_center_semidirectProduct_iff}
\lean{mem_center_semidirectProduct_iff}
\leanfile{M2rGroup7/Lemmas/GroupTheoryLemmas.lean}
\uses{def:fixedPointsSubgroup}
\leanok
Let $N, K$ be abelian groups and $f : K \to \operatorname{Aut}(N)$. For
$g = (n, h) \in N\rtimes_f K$, we have $g \in Z(N\rtimes_f K)$ if and only if
\[
  h \in \ker f
  \qquad\text{and}\qquad
  f(h')(n) = n \text{ for every } h' \in K.
\]
\end{lemma}
\begin{proof}
Write $g = (n, h)$ and let $x = (x', y') \in N\rtimes_f K$ be arbitrary. By
definition of the operation, the commutativity relation $g x = x g$ reads
\[
  \bigl(n\cdot f(h)(x'),\, h y'\bigr)
  \;=\; \bigl(x'\cdot f(y')(n),\, y' h\bigr).
\]
Because $K$ is abelian, the right components automatically agree. Hence
$g\in Z(N\rtimes_f K)$ is equivalent to
\[
  n\cdot f(h)(x') \;=\; x' \cdot f(y')(n)
  \qquad \text{for all } x' \in N,\ y' \in K. \tag{$\ast$}
\]

Setting $y' = 1$ in $(\ast)$ gives $n\cdot f(h)(x') = x' \cdot n$, and since $N$
is abelian we may cancel to get $f(h)(x') = x'$ for every $x' \in N$, i.e.\
$f(h) = \mathrm{id}$, i.e.\ $h \in \ker f$.

Setting $x' = 1$ in $(\ast)$ and using $h \in \ker f$ gives $n = f(y')(n)$ for
every $y' \in K$, i.e.\ $n \in \operatorname{Fix}(\operatorname{Im} f)$.

Conversely, if both conditions hold then $f(h) = \mathrm{id}$ and
$f(y')(n) = n$, so
\[
  (n, h)(x', y') = (n\cdot x',\, h y')
  = (x'\cdot n,\, y' h) = (x', y')(n, h),
\]
using that $N$ and $K$ are abelian.
\end{proof}

\begin{lemma}\label{lem:center_semidirectProduct_iso}
\lean{center_semidirectProduct_iso}
\leanfile{M2rGroup7/Lemmas/GroupTheoryLemmas.lean}
\uses{def:fixedPointsSubgroup, lem:mem_center_semidirectProduct_iff}
\leanok
Let $N, K$ be abelian groups and $f : K \to \operatorname{Aut}(N)$. There is a
group isomorphism
\[
  Z(N\rtimes_f K) \;\simeq\; \operatorname{Fix}(\operatorname{Im} f) \times \ker f
\]
sending the class of $(n, h) \in Z(N\rtimes_f K)$ to $(n, h)$.
\end{lemma}
\begin{proof}
By Lemma~\ref{lem:mem_center_semidirectProduct_iff}, the map
\[
  (n, h) \in Z(N\rtimes_f K) \;\longmapsto\; (n, h) \in
  \operatorname{Fix}(\operatorname{Im} f) \times \ker f
\]
is well defined. It is a homomorphism: for $g_1 = (n_1, h_1)$ and
$g_2 = (n_2, h_2)$ both in the centre, $h_1 \in \ker f$ implies
$f(h_1) = \mathrm{id}$, so
\[
  g_1 g_2 = (n_1 \cdot f(h_1)(n_2),\, h_1 h_2) = (n_1 n_2,\, h_1 h_2),
\]
which matches the componentwise product. Injectivity is immediate, and
surjectivity is the converse direction of
Lemma~\ref{lem:mem_center_semidirectProduct_iff}.
\end{proof}

\end{document}


%%FILE numina/blueprints/factoring-out-lemma-about-center-of-semi-direct-product/source/factoring-out-lemma-about-center-of-semi-direct-product-source.tex
In the proof of this result:

Semidirect products $C_{q^n} \rtimes C_{p^m}$ ($q$ odd prime, $p \neq q$) are classified up to isomorphism by $r \in \{0, …, \min(m, d)\}$ where $d = v_p(q - 1)$, giving $\min(m, d) + 1$ classes. Every action $f$ belongs to exactly one class.

(which has many different forms, one of which is classify_Cqn_rtimes_Cpm, all in CycPGroupClassification.lean)

It uses a lemma that shows that the centre is direct product between Fix(Im(f)) and Ker(f) but I think it only shows for cyclic p and q groups, but its true in general for abelian groups. Could you please take it out and create a lemma that reflects the following?

\begin{lemma}\label{Centre of semidirect product}
    Let $f : K \rightarrow \text{Aut}(N)$ be the underlying homomorphism of $N \rtimes K$ where $N, K$ are abelian groups. Then $Z(N \rtimes K) = \text{Fix}(\text{Im}(f)) \times \text{Ker}(f)$
\end{lemma}
\begin{proof}
    Group operation is defined as $$(x_1, y_1)(x_2, y_2) = (x_1 \cdot f(y_1)(x_2), y_1y_2)$$
    By definition $(x, y) \in Z(G) \iff$ for arbitrary $(x', y') \in G$ we have: $$(x, y)(x', y') = (x', y')(x, y)$$
    Using the definition of the group operation, this is equivalent to $$(x \cdot f(y)(x'), yy') = (x' \cdot f(y')(x), y'y)$$
    Since $K$ is abelian, we always have $yy' = y'y$. Thus the condition simplifies to $$(x, y) \in Z(G) \iff x \cdot f(y)(x') = x' \cdot f(y')(x) \text{ } \forall x' \in N, y' \in K$$
    Setting $y' = 1$ we get $x \cdot f(y)(x') = x' \cdot x \implies f(y)(x') = x'$ (as $N$ is abelian) which shows that $y \in \text{Ker}(f)$ as $x'$ is arbitrary.

    Setting $x' = 1$ we get $x = f(y')(x)$ (using $y \in \text{Ker}(f)$). Hence $x \in \text{Fix}(\text{Im}(f))$ as $y'$ is arbitrary.

    These conditions are necessary and sufficient as one can show that $(x, y)(x', y') = (x \cdot x', y \cdot y') = (x' \cdot x, y' \cdot y) = (x', y') (x, y)$.
\end{proof}

%%FILE numina/blueprints/order-30-classification/order-30-classification.tex
% Blueprint: Order 30 classification
% Generated from the project source LaTeX (Order 30 section).

\begin{document}

\section{Classification of Groups of Order 30}
\label{order-30-classification}

We classify the finite groups of order $30 = 2 \cdot 3 \cdot 5$ up to isomorphism.
The proof follows Deopurkar~\cite{deopurkar_modern_2013}.
There are exactly four isomorphism classes:
\begin{enumerate}
  \item $C_{30}$ (cyclic),
  \item $D_{15}$ (dihedral),
  \item $C_5 \times D_3$ (equivalently $C_5 \times S_3$),
  \item $C_3 \times D_5$.
\end{enumerate}

\subsection{Sylow infrastructure for groups of order \(pqr\)}

\begin{lemma}\label{pqr_normal_sylow}
Let $G$ be a finite group of order $pqr$ for distinct primes $p < q < r$.
Then either the Sylow-$q$ subgroup or the Sylow-$r$ subgroup of $G$ is normal,
i.e.\ $n_q(G) = 1$ or $n_r(G) = 1$.
\end{lemma}

\begin{proof}
By Sylow's theorems $n_r \equiv 1 \pmod r$ and $n_r \mid pq$, so $n_r \in \{1, p, q, pq\}$.
Neither $p$ nor $q$ is congruent to $1$ modulo $r$ because both are smaller than $r$
and at least $2$. If $n_r = 1$ we are done, so assume $n_r = pq$.
Distinct Sylow-$r$ subgroups have prime order $r$, so they intersect trivially;
the total number of elements of order $r$ is therefore $pq(r-1)$.

Now $n_q \mid pr$ and $n_q \equiv 1 \pmod q$. Suppose for contradiction $n_q > 1$.
Then $n_q \geq q + 1$, so $n_q \in \{r, pr\}$ (since $p < q$, $n_q = p$ is impossible).
In either case $n_q \geq r$, giving at least $r(q-1) \geq q+r-1$ elements of order $q$.
Adding the elements of order $r$ and the identity:
\[
1 + pq(r-1) + r(q-1) = 1 + pqr - pq + qr - r > pqr,
\]
since $qr > pq + r - 1$ follows from $r > p$ and $q \geq 3$. Contradiction.
\end{proof}

\begin{theorem}\label{pqr_normal_subgroup_qr}
Let $G$ be a finite group of order $pqr$ for distinct primes $p < q < r$.
Then $G$ has a normal subgroup of order $qr$.
\end{theorem}

\begin{proof}
By Lemma~\ref{pqr_normal_sylow} either $n_q = 1$ or $n_r = 1$.
Pick a Sylow-$q$ subgroup $Q$ and a Sylow-$r$ subgroup $R$.
The two subgroups have coprime orders, so $Q \cap R = \{1\}$ and hence
$|QR| = |Q|\,|R| = qr$.
If $Q$ is normal then $QR$ is a subgroup (as $Q \lhd G$),
of order $qr$ and index $p$. Since $p$ is the smallest prime dividing $|G|$,
the subgroup $QR$ is normal. The case where $R$ is normal is symmetric.
\end{proof}

\subsection{The main classification}

\begin{theorem}\label{order_30_classification}
Let $G$ be a group with $|G| = 30$. Then $G$ is isomorphic to exactly one of
$C_{30}$, $D_{15}$, $C_5 \times D_3$, $C_3 \times D_5$.
\end{theorem}

\begin{proof}
By Theorem~\ref{pqr_normal_subgroup_qr} (with $p=2$, $q=3$, $r=5$), $G$ has a
normal subgroup $N$ of order $15$. By Schur--Zassenhaus there is a complementary
subgroup $U \leq G$ of order $2$ with $G \cong N \rtimes U$.
We have $U \cong C_2$, and by the classification of groups of order $pq$ with
$p \nmid q-1$ (taking $p=3$, $q=5$, where $3 \nmid 4$) we have $N \cong C_{15}$.

Let $\phi : C_2 \to \mathrm{Aut}(C_{15})$ be the underlying homomorphism.
The image $\mathrm{Im}\,\phi$ has order $1$ or $2$.

\paragraph{Case 1: $|\mathrm{Im}\,\phi| = 1$.}
The action is trivial, so $G \cong C_{15} \times C_2 \cong C_{30}$.

\paragraph{Case 2: $|\mathrm{Im}\,\phi| = 2$.}
We have $\mathrm{Aut}(C_{15}) \cong (\mathbb{Z}/15\mathbb{Z})^\times
\cong (\mathbb{Z}/3)^\times \times (\mathbb{Z}/5)^\times \cong C_2 \times C_4$,
which has exactly three elements of order two: those acting by
\begin{itemize}
  \item inverting both factors (the global inversion automorphism of $C_{15}$),
  \item inverting only the $C_3$ factor,
  \item inverting only the $C_5$ factor.
\end{itemize}
By Lemma~\ref{equal_image_iso_sdp} (\textit{equal images give isomorphic
semi-direct products for cyclic $p$-groups in the acting position}), each choice
gives one isomorphism class. The three resulting groups, written in their direct
product form, are:
\begin{itemize}
  \item $\phi$ inverts both $C_3$ and $C_5$: $G \cong D_{15}$.
  \item $\phi$ inverts only $C_3$ and fixes $C_5$:
        $G \cong (C_3 \rtimes_{\mathrm{inv}} C_2) \times C_5 \cong D_3 \times C_5$.
  \item $\phi$ fixes $C_3$ and inverts $C_5$:
        $G \cong C_3 \times (C_5 \rtimes_{\mathrm{inv}} C_2) \cong C_3 \times D_5$.
\end{itemize}
Combined with Case~1, we obtain the four isomorphism classes claimed.
\end{proof}

\begin{lemma}\label{equal_image_iso_sdp}
If $K$ is a cyclic $p$-group, two homomorphisms $f, g : K \to \mathrm{Aut}(H)$
with the same image define isomorphic semi-direct products $H \rtimes_f K$ and
$H \rtimes_g K$.
\end{lemma}

\begin{proof}
Let $k$ be a generator of $K$. Since $\mathrm{Im}(f) = \langle f(k) \rangle =
\langle g(k) \rangle$, there exists $n$ with $\gcd(n, |\mathrm{Im}(f)|) = 1$
such that $g(k) = f(k)^n$. As $\mathrm{Im}(f)$ has order $p^j$ for some $j \geq 0$,
this $n$ is coprime to $|K| = p^m$. The map $\beta(x) = x^n$ is then an
automorphism of $K$, and $g = f \circ \beta$. The pair $(\mathrm{id}_H, \beta)$
induces an isomorphism $H \rtimes_g K \cong H \rtimes_f K$.
\end{proof}

\subsection{Uniqueness}

The four isomorphism classes are pairwise non-isomorphic. This is verified by
the decidable group invariants framework
(\verb|M2rGroup7.Uniqueness.uniqueness|), which uses a computable family of
invariants (order, abelianness, element count by order, etc.) and checks all
relevant pairs via \verb|native_decide|. In particular:
\begin{itemize}
  \item $C_{30}$ is the unique abelian group of order $30$.
  \item $D_{15}$ has $15$ elements of order $2$, while $C_5 \times D_3$ has
        $3$ and $C_3 \times D_5$ has $5$ elements of order $2$.
\end{itemize}
This distinguishes all four classes.

\end{document}


%%FILE numina/blueprints/order-30-classification/source/order-30-classification-source.tex
Follow the Order 30 section of this document:

%%%%%%%%%%%%%%%%%%%%%%%%%%%%%%%%%%%%%%%%%%
% Mathematics Final Year Research Projects
% LaTeX Template
% Version 1.0 (31/01/24)
%
% This template has been adapted from: https://www.overleaf.com/latex/templates/imperial-college-report-template/wncnzptkhnbc
% Students should feel free to adapt this template to their needs.
%%%%%%%%%%%%%%%%%%%%%%%%%%%%%%%%%%%%%%%%%%
%----------------------------------------------------------------------------------------
% PACKAGES AND OTHER DOCUMENT CONFIGURATIONS
%----------------------------------------------------------------------------------------
\documentclass[a4paper,11pt, titlepage]{article}

%% Language and font encodings
\usepackage[english]{babel}
\usepackage[utf8]{inputenc}
\usepackage[T1]{fontenc}

%% Sets page size and margins
\usepackage[a4paper,top=1in,bottom=1in,left=1in,right=1in,marginparwidth=1.75cm]{geometry}

\usepackage{xcolor}

%% Useful packages
\usepackage{afterpage}
\usepackage{amsmath}
\usepackage{amsthm}
\usepackage{amssymb}
\usepackage{thmtools}
\usepackage{mdframed}
\usepackage{csquotes}
\usepackage{enumitem}
\usepackage{graphicx}
\usepackage{lipsum}
\usepackage{booktabs}
\usepackage{unicode-math}
\usepackage[table]{xcolor} % Required for \cellcolor
\usepackage{microtype}
\usepackage{fontspec}
\usepackage{minted}
\usepackage{titlesec}

% Define a soft, professional green for cell backgrounds
\definecolor{softgreen}{HTML}{E2F0D9} 

% Switch to a monospace font supporting Unicode characters
\newfontfamily{\leanCodeFont}{JuliaMono}[NFSSFamily=leanfont]

\mdfdefinestyle{fancycode}{
  linewidth=2pt,
  rightline=false, topline=false, bottomline=false,
  linecolor=gray!60!black,       % A neutral dark gray bar for code
  backgroundcolor=gray!5!white,  % Subtle light gray background
  innerbottommargin=8pt,
  innertopmargin=8pt,
  leftmargin=1.5em,              % Matches the indent of your Definitions
  rightmargin=1.5em,
}

\usemintedstyle{tango}
\newmintinline[lean]{lean4}{fontfamily=leanfont, bgcolor=gray!5!white}
\newminted[leancode]{lean4}{fontfamily=leanfont, fontsize=\footnotesize, breaklines=true}
\surroundwithmdframed[style=fancycode]{leancode}


% Listings (for displaying code):
\usepackage{listings}
\lstset{
    frame = single, 
    framexleftmargin=15pt
}

% Center figure captions:
\usepackage{caption}
\captionsetup[figure]{labelfont={bf},name={Figure},labelsep=quad}
\captionsetup[table]{labelfont={bf},name={Table},labelsep=quad}


% ----------- Algorithm2e setup
\usepackage[ruled,vlined]{algorithm2e}
\makeatletter
\renewcommand{\SetKwInOut}[2]{%
  \sbox\algocf@inoutbox{\KwSty{#2}\algocf@typo:}%
  \expandafter\ifx\csname InOutSizeDefined\endcsname\relax% if first time used
    \newcommand\InOutSizeDefined{}\setlength{\inoutsize}{\wd\algocf@inoutbox}%
    \sbox\algocf@inoutbox{\parbox[t]{\inoutsize}{\KwSty{#2}\algocf@typo:\hfill}~}\setlength{\inoutindent}{\wd\algocf@inoutbox}%
  \else% else keep the larger dimension
    \ifdim\wd\algocf@inoutbox>\inoutsize%
    \setlength{\inoutsize}{\wd\algocf@inoutbox}%
    \sbox\algocf@inoutbox{\parbox[t]{\inoutsize}{\KwSty{#2}\algocf@typo:\hfill}~}\setlength{\inoutindent}{\wd\algocf@inoutbox}%
    \fi%
  \fi% the dimension of the box is now defined.
  \algocf@newcommand{#1}[1]{%
    \ifthenelse{\boolean{algocf@inoutnumbered}}{\relax}{\everypar={\relax}}%
%     {\let\\\algocf@newinout\hangindent=\wd\algocf@inoutbox\hangafter=1\parbox[t]{\inoutsize}{\KwSty{#2}\algocf@typo\hfill:}~##1\par}%
    {\let\\\algocf@newinout\hangindent=\inoutindent\hangafter=1\parbox[t]{\inoutsize}{\KwSty{#2}\algocf@typo:\hfill}~##1\par}%
    \algocf@linesnumbered% reset the numbering of the lines
  }}%
\makeatother
% --------- end algorithm2e setup

% \bm allows typing bold math:
\usepackage{bm}
\usepackage[normalem]{ulem}

\usepackage[colorinlistoftodos]{todonotes}
\usepackage[colorlinks=true, allcolors=blue]{hyperref}

\renewcommand*{\rmdefault}{bch}
\renewcommand*{\ttdefault}{lmtt}
\newcommand{\citationneeded}{\textcolor{red}{[citation-needed]}}

\DeclareMathOperator*{\argmin}{\arg\!\min}
\DeclareMathOperator*{\argmax}{\arg\!\max}

% Add bigger skip between paragraphs, makes reading easier:
\setlength{\parskip}{0.5em}

% Bibliography
\usepackage[numbers, comma, square, sort&compress]{natbib}
\bibliographystyle{abbrvunsrtnat.bst}


% Theorem, Lemma, Corollary and Proofs
% \theoremstyle{definition}
% \newtheorem{definition}{Definition}[section]
% \theoremstyle{plain}
% \newtheorem{theorem}{Theorem}[section]
% \newtheorem{corollary}{Corollary}[theorem]
% \newtheorem{lemma}[theorem]{Lemma}
% \theoremstyle{remark}
% \newtheorem*{remark}{Remark}
% --- Define the style ---
\declaretheoremstyle[
  headfont=\bfseries\color{blue!70!black},
  bodyfont=\itshape,
  mdframed={
    linewidth=2pt,
    rightline=false, topline=false, bottomline=false, % Only left line
    linecolor=blue!70!black,
    backgroundcolor=blue!5!white,
    innerbottommargin=8pt,
    innertopmargin=8pt,
  }
]{fancytheorem}

\declaretheoremstyle[
  headfont=\bfseries\color{green!50!black},
  bodyfont=\normalfont,
  mdframed={
    linewidth=2pt,
    rightline=false, topline=false, bottomline=false,
    linecolor=green!50!black,
    backgroundcolor=green!5!white,
  }
]{fancydef}

% --- Apply the styles ---
\declaretheorem[style=fancytheorem, name=Theorem, numberwithin=section]{theorem}
\declaretheorem[style=fancytheorem, name=Lemma, sibling=theorem]{lemma}
\declaretheorem[style=fancytheorem, name=Corollary, sibling=theorem]{corollary}
\declaretheorem[style=fancydef, name=Definition, numberwithin=section]{definition}

% --- Keep Remark Simple ---
\theoremstyle{remark}
\newtheorem*{remark}{Remark}

\titlespacing*{\section}{0pt}{3.5ex plus 1ex minus .2ex}{2ex plus .2ex}
\titlespacing*{\subsection}{0pt}{2.5ex plus 1ex minus .2ex}{1ex plus .2ex}

% --- Section: Blue with matching underline ---
\titleformat{\section}
  {\normalfont\Large\bfseries\color{black}} % Font and color
  {\thesection}                                     % The section number
  {1em}                                             % Space between number and text
  {}                                                % Code applied to the title text
  [\vspace{2pt}\color{blue!70!black}\titlerule]     % The line underneath

% --- Subsection: Matching color, no line ---
\titleformat{\subsection}
  {\normalfont\large\bfseries\color{black}}
  {\thesubsection}
  {1em}
  {}

% paragraph headings
\setcounter{secnumdepth}{4}

% Custom symbols
\newcommand{\isom}{\stackrel{\sim}{\smash{\rightarrow}\rule{0pt}{0.8ex}}}

%----------------------------------------------------------------------------------------
% END OF DOCUMENT CONFIGURATION
%----------------------------------------------------------------------------------------

%----------------------------------------------------------------------------------------
% IMPORTANT INFORMATION TO MODIFY FOR THE TITLE PAGE
%----------------------------------------------------------------------------------------

% Kevin: we should maybe call this formalising a databse of small groups
\newcommand{\reporttitle}{Formalising Groups of Small Order}
\newcommand{\reportauthorA}{Sevin Fernando (CID: 02568398)}
\newcommand{\reportauthorD}{Jay Lees-Maffei (CID: 02595674)}
\newcommand{\reportauthorC}{Tianyu Qi (CID: 02556308)}
\newcommand{\reportauthorB}{Agamjeet Singh (CID: 02593148)}
\newcommand{\supervisor}{Kevin Buzzard}
%----------------------------------------------------------------------------------------

%----------------------------------------------------------------------------------------
% To compile this file, use the following sequence: 
%	latex main.tex
%	bibtex main.tex
%	latex main.tex
%	latex main.tex
%----------------------------------------------------------------------------------------

%----------------------------------------------------------------------------------------
% START OF DOCUMENT
%----------------------------------------------------------------------------------------
\begin{document}

%%%%%%%%%%%%%%%%%%%%%%%%%%%%%%%%%%%%%
% Title page
%%%%%%%%%%%%%%%%%%%%%%%%%%%%%%%%%%%%%
\begin{titlepage}
\newcommand{\HRule}{\rule{\linewidth}{0.5mm}} % Defines a new command for the horizontal lines, change thickness here
%----------------------------------------------------------------------------------------
%	LOGO SECTION
%----------------------------------------------------------------------------------------
\includegraphics[width=8cm]{Imperial_logo.png}\\[1cm] % Include a department/university logo - this will require the graphicx package
%----------------------------------------------------------------------------------------
\center % Center everything on the page
%----------------------------------------------------------------------------------------
%	HEADING SECTIONS
%----------------------------------------------------------------------------------------
% For M3R or M4R reports, comment out as appropriate
\textsc{\LARGE Imperial College London}\\[0.5cm] 
\textsc{\Large Department of Mathematics}\\[1.5cm] 
\textsc{\Large Second-year Group Research Project}\\[0.5cm]
%----------------------------------------------------------------------------------------
%	TITLE SECTION
%----------------------------------------------------------------------------------------
\makeatletter
\HRule \\[0.6cm]
{ \huge \bfseries \reporttitle}\\[0.6cm] % Title of your document
\HRule \\[1.5cm]
%----------------------------------------------------------------------------------------
%	AUTHOR SECTION
%----------------------------------------------------------------------------------------
\begin{minipage}{0.4\textwidth}
\begin{flushleft} \large
\emph{Author:}\\
\reportauthorA \\
\reportauthorB \\
\reportauthorC \\
\reportauthorD
\end{flushleft}
\end{minipage}
~
\begin{minipage}{0.4\textwidth}
\begin{flushright} \large
\emph{Supervisor:} \\
\supervisor % Supervisor's name
\end{flushright}
\end{minipage}\\[2cm]
\makeatother
%----------------------------------------------------------------------------------------
%	FOOTER & DATE SECTION
%----------------------------------------------------------------------------------------
\vfill
\makeatletter
{\large \today}\\[2cm] % Date, change the \today to a set date if you want to be precise
\makeatother
\end{titlepage}
%%%%%%%%%%%%%%%%%%%%%%%%%%%%%%%%%%%%%

%%%%%%%%%%%%%%%%%%%%%%%%%%%%%%%%%%%%%
% Abstract
\begin{abstract}
Computer-assisted formalisation of mathematics is a highly active field of research which focuses on translating mathematical statements into computer verifiable code in order to ensure correctness, conciseness, and allows the possibility for new proofs to be discovered manually or automatically. On the other hand, classifying finite groups up to isomorphism is an important result in group theory that, while heavily cross-checked amongst computer algebra systems such as GAP or MAGMA, has not been formally proven.
This project aims to reduce the gap by formalising the classification of finite groups of small order in the Lean 4 programming language.
To achieve this goal, we constructed a library of explicit groups which represent the isomorphism classes. We adapted/translated existing proofs into two main theorems to prove completeness and uniqueness. In doing this, we employed novel methodology and tooling such as Claude Code, Numina Fuse, and Google Gemini to leverage the mathematical reasoning and auto-proving capability of modern large language models.
[WIP]
\end{abstract}
%%%%%%%%%%%%%%%%%%%%%%%%%%%%%%%%%%%%%
 
%%%%%%%%%%%%%%%%%%%%%%%%%%%%%%%%%%%%%
% Table of contents
\tableofcontents
\clearpage
%%%%%%%%%%%%%%%%%%%%%%%%%%%%%%%%%%%%%

%%%%%%%%%%%%%%%%%%%%%%%%%%%%%%%%%%%%%
% Main sections
%%%%%%%%%%%%%%%%%%%%%%%%%%%%%%%%%%%%%
\section{Introduction}
\label{sec:introduction}


%The introduction should attempt to set your work in the context of other work done in the field. It
%should demonstrate that you are aware of what you are doing, and how it relates to other work
%(with references). It should also provide an overview of the contents of the project. You should
%highlight your individual contributions and any novel result: which of the calculations, theorems,
%examples, proofs, conjectures, codes etc. are your own?



% Kevin: we are not formalising the GAP small groups library, this is too big
In this project we aim to formalise the GAP small groups library \cite{noauthor_gap_nodate}.
This involves implementing the groups in a proof assistant and proving two theorems:
\begin{itemize}
\item Any group of small enough order is isomorphic to a group in the library (classification).
\item No two groups in the library are isomorphic (uniqueness).
\end{itemize}

To do this we will use Lean, a proof assistant and dependently typed programming language \cite{noauthor_lean_nodate}.

There have been other proofs classifying groups of small order in Lean, but none operating
on a library of concrete groups.
(Of particular use in our project is Harper and Wu's classification of groups of order \(pq\)
for prime \(p, q\) \cite{harper_classifying_2025}.)
However, most proofs are missing from Lean's Mathlib \cite{noauthor_mathlib4_nodate}
and we have had to formalise these ourselves.
This includes classification theorems for groups of order \(p^2q\) and \(p^3\) for prime \(p, q\).
These newly formalised classification theorems make up the bulk of our project and are results
that we hope will be of use to other group theorists interested in formalisation.

Other important results from our project include the formalised classification theorems
for groups of order \(16\) and \(24\).
These are not general cases, so we do not expect them to be of much interest to other mathematicians,
but they are still necessary part of our overall classification theorem. 

\begin{table}[htbp]
\centering
\caption{Classification of Groups of Order 1 to 31}
\label{tab:groups_1_31_clean}
\small
\begin{tabular}{ccc | ccc | ccc}
\hline
Order & Pattern & Reference & Order & Pattern & Reference & Order & Pattern & Reference \\
\hline
1  & —      & —                                      & 12 & $p^2 q$ & \ref{$p^2 q$ for distinct primes $p, q$} & 23 & $p$     & \ref{$p$ for prime $p$} \\
2  & $p$    & \ref{$p$ for prime $p$}                & 13 & $p$     & \ref{$p$ for prime $p$}                & 24 & $p^3 q$ & \ref{Order $24$} \\
3  & $p$    & \ref{$p$ for prime $p$}                & 14 & $pq$    & \ref{$pq$ for distinct primes $p, q$}  & 25 & $p^2$   & \ref{$p^2$ for prime $p$} \\
4  & $p^2$  & \ref{$p^2$ for prime $p$}              & 15 & $pq$    & \ref{$pq$ for distinct primes $p, q$}  & 26 & $pq$    & \ref{$pq$ for distinct primes $p, q$} \\
5  & $p$    & \ref{$p$ for prime $p$}                & 16 & $p^4$   & \ref{Order $16$}                       & 27 & $p^3$   & \ref{$p^3$ for prime $p$} \\
6  & $pq$   & \ref{$pq$ for distinct primes $p, q$}  & 17 & $p$     & \ref{$p$ for prime $p$}                & 28 & $p^2 q$ & \ref{$p^2 q$ for distinct primes $p, q$} \\
7  & $p$    & \ref{$p$ for prime $p$}                & 18 & $p^2 q$ & \ref{$p^2 q$ for distinct primes $p, q$} & 29 & $p$     & \ref{$p$ for prime $p$} \\
8  & $p^3$  & \ref{$p^3$ for prime $p$}              & 19 & $p$     & \ref{$p$ for prime $p$}                & 30 & $p q r$ & \ref{Order $30$} \\
9  & $p^2$  & \ref{$p^2$ for prime $p$}              & 20 & $p^2 q$ & \ref{$p^2 q$ for distinct primes $p, q$} & 31 & $p$     & \ref{$p$ for prime $p$} \\
10 & $pq$   & \ref{$pq$ for distinct primes $p, q$}  & 21 & $pq$    & \ref{$pq$ for distinct primes $p, q$}  &    &         & \\
11 & $p$    & \ref{$p$ for prime $p$}                & 22 & $pq$    & \ref{$pq$ for distinct primes $p, q$}  &    &         & \\
\hline
\end{tabular}
\end{table}

\subsection{Lean}
We wrote earlier that Lean is dependently typed.

\subsection{Semidirect Product}

\subsection{Sylow's Theorem}

We extensively employ Sylow's theorems for the classification of groups of orders $pq, p^2q, 24$ and $30$. A standard notation we use frequently in this report is that $v_p(n)$ denotes the highest power of a prime $p$ dividing $n$.

We state the results and their consequences briefly:

\begin{theorem}[First Sylow Theorem]
    Let $G$ be a finite group and $p$ be a prime such that $p \mid |G|$. Then $G$ has a subgroup of order $p^{v_p(|G|)}$.
\end{theorem}

We would use the standard notation that $n_p(G)$ represents the number of Sylow-$p$ groups of $G$.

\begin{theorem}[Second Sylow Theorem]
    Let $G$ be a finite group and $p$ be a prime such that $p \mid |G|$. Then $n_p(G) \equiv 1 \pmod p$.
\end{theorem}
\begin{theorem}[Fourth Sylow Theorem]
    Let $G$ be a finite group and $p$ be a prime such that $p \mid |G|$. The Sylow-$p$ groups of $G$ form a single conjugacy class.
\end{theorem}
A very useful corollary of the above theorem is (using Orbit-Stabiliser theorem):
\begin{corollary}
    Let $G$ be a finite group and $p$ be a prime such that $p \mid |G|$. Then $n_p(G) \mid |G|$.
\end{corollary}

A further corollary is that $n_p(G) \mid \frac{|G|}{p^{v_p(|G|)}}$ as $n_p(G) \equiv 1 \pmod p$ implies $p \nmid n_p(G)$.

If the group $G$ is clear from context, we would often abbreviate $n_p(G)$ as just $n_p$.

The constraints $n_p \equiv 1 \pmod p$ and $n_p \mid |G|$ in particular are extremely useful and will be employed frequently.

In conjunction with these, we would use the following corollary often too (follows from Fourth Sylow Theorem and definition of normality)
\begin{corollary}
    Let $G$ be a finite group. If $n_p(G) = 1$, then the unique Sylow-$p$ group of $G$ is normal.
\end{corollary}

\subsection{Schur-zassenhaus's theorem}

The Schur-zassenhaus's theorem is useful in showing that a group $G$ must be a semi-direct product when given a normal subgroup $N$.

\begin{theorem}[Schur-zassenhaus]
    Let $G$ be a finite group and let $N$ be a normal subgroup. Further suppose that $|N|$ and $|G/N|$ are co-prime. Then $G$ is isomorphic to some semi-direct product $N \rtimes G/N$.
\end{theorem}

This theorem is extremely useful in the proofs for groups of order $pq, p^2 q$ and $pqr$ because all such groups are semi-direct products.

\section{Classification}

The classification of groups of a particular order depends on the shape of its prime factorisation. So we divided the classification into the following sections:

\subsection{$p$ for prime $p$}
\label{$p$ for prime $p$}

\subsection{$p^2$ for prime $p$}
\label{$p^2$ for prime $p$}

\subsection{$pq$ for distinct primes $p, q$}
\label{$pq$ for distinct primes $p, q$}

\subsection{$p^3$ for prime $p$}
\label{$p^3$ for prime $p$}

\subsection{$p^2 q$ for distinct primes $p, q$}
\label{$p^2 q$ for distinct primes $p, q$}

Classification of groups of order $p^2 q$ can be done solely using semi-direct products because of this result:
\begin{lemma}[With Claude]\label{Groups of order $p^2 q$ have normal Sylow subgroup}
    All groups $G$ of order $p^2 q$ have either a normal Sylow-$p$ group or a normal Sylow $q$ group.
\end{lemma}
\begin{proof}
    By Sylow's theorems, $n_p \mid q$, $n_p \equiv 1 \pmod p$, $n_q \mid p^2$, $n_q \equiv 1 \pmod q$. Hence $n_q$ is either $1, p$ or $p^2$ and $n_p$ is either $1$ or $q$.
    
    Case 1: $p < q$
    
    If $n_q = p$ then $p \equiv 1 \pmod q$, contradicting $p < q$ as $p \neq 1$.
    If $n_q = p^2$ then we can perform a counting argument. We have $p^2$ Sylow-$q$ groups and they are trivially intersecting as they have prime order. Hence we get $p^2 (q - 1)$ elements of order $q$. As $n_p \geq 1$, the remaining $p^2$ elements must form a Sylow-$p$ group, hence $n_p = 1$, and so there is a normal Sylow-$p$ group.

    Otherwise $n_q = 1$ and so there exists a normal Sylow-$q$ group in that case.

    Case 2: $p > q$
    
    If $n_p = q$, then we have a contradiction to $n_p \equiv 1 \pmod p$ as $q < p$ and $q \neq 1$. Hence $n_p = 1$ and there exists a normal Sylow-$p$ group. 
\end{proof}

\begin{lemma}[With Gemini]\label{$K$ cyclic $p$-group then equal images define isomorphic semi-direct products}
    If $K$ is a cyclic $p$-group, two homomorphisms $f, g : K \rightarrow \text{Aut}(H)$ define isomorphic semi-direct products $H \rtimes K$ if they have the same image.
\end{lemma}
\begin{proof}
    It suffices to find an automorphism $\beta \in \text{Aut}(H)$ such that $g = f \circ \beta$. Let $k$ be a generator of $K$. We have that $\text{Im}(f) = \langle f(k) \rangle$ and $\text{Im}(g) = \langle g(k) \rangle$. As the images are equal, we get that $g(k) \in \langle f(k) \rangle$ and so there exists $n$ such that $g(k) = f(k)^n$ where $\gcd(n, |\text{Im}(f)|) = 1$. By the first isomorphism theorem, $\text{Im}(f) \cong K / \ker(f)$, hence $\text{Im}(f)$ has order $p^j$ for $0 \leq j \leq m$. 
    
    If $j = 0$ then we get that the image of $f, g$ is the trivial group, hence both of them lead to direct products which are clearly isomorphic. 
    
    Otherwise $j > 0$ and we get that $\gcd(n, p^j) = 1$ so $p \nmid n$. Hence $\gcd(n, |K|) = 1$. So we can define the automorphism $\beta : K \rightarrow K$ as $\beta(x) = x^n$. It is well defined as $(xy)^n = x^n y^n$ follows from the fact that $K$ is abelian. It is injective because $x^n = 1$ implies $x = 1$ as the order of $x$ would have to divide $n$ and $|K|$ (by Lagrange's theorem) but their gcd is $1$. Surjectivity then follows from the fact that it is injective and that it is a map from a finite group to itself.

    We have $g = f \circ \beta$ as they agree on a generator since $f(\beta(k)) = f(k^n) = f(k)^n = g(k)$.
\end{proof}

\begin{theorem}[With Gemini and Numina Fuse]\label{Classification of semi-direct products between Cyclic p-group and q-groups}
    Semi-direct products $C_{q^n} \rtimes C_{p^m}$ ($q$ odd prime, $p \neq q$) are classified up to isomorphism by $r \in \{0, …, \min(m, d)\}$ where $d = v_p(q - 1)$, giving $\min(m, d) + 1$ classes. Every action $f$ belongs to exactly one class.
\end{theorem}
\begin{proof}
    The fact that every action belongs to at least one class is a lot easier than showing that it belongs to exactly one class.

    We know that the semi direct product must have a homomorphism $f : C_{p^m} \rightarrow \text{Aut}(C_{q^n})$. We have $\text{Aut}(C_{q^n}) \cong C_{q^{n-1}(q - 1)}$. Thus, the image of $f$ (by first isomorphism theorem) can have order $p^r$ for some $0 \leq r \leq \min(v_p(q-1), m)$. Then we are done by the fact that a cyclic group has a unique subgroup given an order and that equal images imply isomorphic semi direct products (as $C_{p^m}$ is a cyclic $p$-group). For each $0 \leq r \leq \min(v_p(q-1), m)$ there does exist a homomorphism with image of order $p^r$. This can be constructed by sending the generator of $C_{p^m}$ to the element of $C_{q^{n-1}(q-1)}$ of order $p^r$.

    To show uniqueness, we will compute the orders of the centres given $r$.

    Claim 1: $Z(G) = \text{Fix}(\text{Im}(f)) \times \text{Ker}(f)$

    Proof: Group operation is defined as $(x_1, y_1)(x_2, y_2) = (x_1 \cdot f(y_1)(x_2), y_1y_2)$. Thus $(x, y) \in Z(G) \iff$ for arbitrary $(x', y') \in G$ we have: $(x, y)(x', y') = (x', y')(x, y)$ or $(x \cdot f(y)(x'), yy') = (x' \cdot f(y')(x), y'y)$. As $C_{p^m}$ is abelian, we always have $yy' = y'y$. Thus $(x, y) \in Z(G) \iff x \cdot f(y)(x') = x' \cdot f(y')(x) \forall x' \in C_{q^n}, y' \in C_{p^m}$. 
    
    Setting $y' = 1$ we get $x \cdot f(y)(x') = x' \cdot x \implies f(y)(x') = x'$ (as $C_{q^n}$ is abelian) which shows that $y \in \text{Ker}(f)$ as $x'$ is arbitrary.

    Setting $x' = 1$ we get $x = f(y')(x)$ (using $y \in \text{Ker}(f)$). Hence $x \in \text{Fix}(\text{Im}(f))$ as $y'$ is arbitrary.

    These conditions are necessary and sufficient as one can show that $(x, y)(x', y') = (x \cdot x', y \cdot y') = (x' \cdot x, y' \cdot y) = (x', y') (x, y)$. $\blacksquare$

    If $r = 0$, the action is trivial, $\text{Im}(f) = \{\text{id}\}$ so every point gets fixed $\text{Fix}(\text{Im}(f)) = C_{q^n}$. Also $\text{Ker}(f) = C_{p^m}$ so $|Z(G)| = p^m q^n$.

    If $r > 0$, we have $|\text{Im}(f)| = p^r$. By the first Isomorphism theorem, $|\text{Ker}(f)| = \frac{|C_{p^m}|}{|\text{Im}(f)|} = p^{m - r}$.

    Claim 2: If $r > 0$, $|\text{Fix}(\text{Im}(f))| = 1$

    Proof: A well known lemma (\mintinline{lean4}{IsPGroup.smul_mul_inv_trivial_or_surjective} in Mathlib) states that if a group $K$ acts on a cyclic $q$-group $G$ of coprime order, the map $\Psi : K \times G \to G$ defined by $(k, g) \mapsto k \bullet g \cdot g^{-1}$ is either trivial ($\Psi(k, g) = 1 \forall k, g \in K \times G$) or surjective. 
    
    In our case, it is not surjective as that would mean $f$ has a trivial image but the image has order $p^r$ with $r > 0$. Hence $\Psi$ must be surjective. 
    
    Let $x \in C_{q^n}$ be an arbitrary fixed point of the action. It suffices to show that $x = 1$. Using surjectivity of $\Psi$, there exists $k_0 \in C_{p^m}$ and $q_0 \in C_{q^n}$ satisfty $\Psi(k_0, q_0) = x \implies k_0 \bullet q_0 = x \cdot q_0$. 
    
    One can inductively show that $k_0^s \bullet (k_0 \bullet q_0) = x^s \cdot (k_0 \bullet q_0)$ for any $s \geq 0$ (Sketch: $k_0 \bullet (k_0 \bullet q_0) = k_0 \bullet (x \cdot q_0) = (k_0 \bullet x) \cdot (k_0 \bullet q_0) = x \cdot (k_0 \bullet q_0)$). 
    
    Substituting $s = p^m$ gives $k_0 \bullet q_0 = x^{p^m} \cdot (k_0 \bullet q_0)$ as $k_0$ has order $p^m$ because $k_0 \in C_{p^m}$. Hence $x^{p^m} = 1$. But $x \in C_{q^n}$ hence $x = 1$ as $p \neq q$. $\blacksquare$

    Hence $|Z(G)| = p^{m - r}$.

    Clearly the order of all centres are distinct, which shows uniqueness.
\end{proof}

Let $s_m(p)$ represent the number of conjugacy classes of matrices in $GL_2(p)$ of order $m$.

\begin{lemma}[Numina Fuse]
    $s_2(p) = 2$ if $p$ is an odd prime.
\end{lemma}
\begin{proof}
    We give both a short and elegant proof and a bashy one (by Numina Fuse) that is much simpler to formalise.

    Let $A \in GL_2(p)$ be of order $2$. Then its minimal polynomial divides $x^2 - 1$, hence it is a product of distinct linear factors ($1 \neq -1$ as $p$ is odd) and so is diagonalisable. The only matrix similar to identity is itself (and it doesn't have order $2$). Hence, $A$ must be similar to one of $\text{diag}(-1, -1)$, $\text{diag}(-1, 1)$ ($\text{diag}(1, -1)$ is similar to $\text{diag}(-1, 1)$) giving two similarity or conjugacy classes.

    The proof we used for formalisation is different and was generated  by Numina Fuse:
    [VERIFY]
    
    Assume $A = \begin{pmatrix} a & b \\ c & d \end{pmatrix}$ satisfies $A^2 = I$ and $A \neq I$. If $A = -I$, we are done. Otherwise, assume $A \neq \pm I$. $A^2 = I$ implies: $a^2 + bc = 1$, $c(a + d) = 0$, $b(a + d) = 0$ and $bc + d^2 = 1$. 

    If $c \neq 0$, the equation $c(a + d) = 0$ forces $d = -a$. The matrix $P = \begin{pmatrix} a+1 & a-1 \\ c & c \end{pmatrix}$ satisfies $P^{-1}AP = \text{diag}(1, -1)$.

    If $c = 0$, the equations reduce to $a^2 = 1$ and $d^2 = 1$, so $a, d \in \{1, -1\}$. Since $A \neq I$ and $A \neq -I$, we must have $a \neq d$. This leaves two possible upper triangular matrices: $\begin{pmatrix} 1 & b \\ 0 & -1 \end{pmatrix}$ and $\begin{pmatrix} -1 & b \\ 0 & 1 \end{pmatrix}$. Both are conjugate to $\text{diag}(1, -1)$.
\end{proof}

\begin{lemma}
    $s_q(p) = \frac{q+3}{2} \Delta^q_{p-1} + \Delta^q_{p+1}$ if $q$ is an odd prime and $p \neq q$.
\end{lemma}
\begin{proof}
    TODO
\end{proof}

\begin{lemma}
    Any two non-trivial semi-direct products $C_q \rtimes (C_p \times C_p)$ arising from homomorphisms with images of order $p$ are isomorphic.
\end{lemma}
\begin{proof}
    Let $f, g : C_p \times C_p \rightarrow C_{q-1}$ be the underlying homomorphisms. It suffices to find an automorphism $\beta \in \text{Aut}(H)$ such that $g = f \circ \beta$. 
\end{proof}

\subsubsection{Counting arguments in Lean}

It is non-trivial to perform counting arguments in lean. We divide it into the following lemmas:

\begin{leancode}
/-- Two distinct Sylow p-subgroups, each of prime order p, intersect trivially. -/
lemma sylow_prime_order_disjoint {p : ℕ} [hp : Fact p.Prime]
    {G : Type*} [Group G] [Finite G]
    (h_sylow_card : ∀ P : Sylow p G, Nat.card ↥(P : Subgroup G) = p)
    {P₁ P₂ : Sylow p G} (hne : P₁ ≠ P₂) :
    Disjoint (P₁ : Subgroup G) P₂
\end{leancode}

\begin{leancode}
/-- The union of all Sylow p-subgroups, when each has prime order p, has
    cardinality 1 + n_p * (p - 1) where n_p = Nat.card (Sylow p G). -/
lemma sylow_prime_union_card {p : ℕ} [hp : Fact p.Prime]
    {G : Type*} [Group G] [Finite G]
    (h_sylow_card : ∀ P : Sylow p G, Nat.card ↥(P : Subgroup G) = p) :
    Nat.card (⋃ P : Sylow p G, (P : Subgroup G) : Set G) =
      1 + Nat.card (Sylow p G) * (p - 1)
\end{leancode}

We use these to show that $G \backslash U$ has size $p^2 - 1$, where $U$ is the union of all Sylow-$q$ subgroups. Every Sylow-$p$ subgroup has $p^2 - 1$ non-identity elements, all outside $U$. By size, each Sylow-$p$ subgroup's non-identity elements equal $G \backslash U$ exactly $\implies$ all Sylow-$p$ subgroups are equal $\implies n_p = 1$.

This infrastructure would come in handy for proving that any group of order $30$ has a normal Sylow-$5$ subgroup.

\subsubsection{The main proof}

We know that there exists either a normal Sylow-$p$ group or a normal Sylow-$q$ group by \ref{Groups of order $p^2 q$ have normal Sylow subgroup}.

\paragraph{Case 1:} There exists a normal Sylow $p$-group

Let $N$ be the normal Sylow $p$-group. Then $|N| = p^2$ and by Schur-zassenhaus theorem there exists a complementary subgroup $U$ of $G$ with $|U| = q$ such that $G \cong N \rtimes U$. We have $U \cong C_q$ and either $N \cong C_{p^2}$ or $N \cong C_p \times C_p$ (by classification of groups of order $p^2$, TODO - link result).

If $N \cong C_{p^2}$, then we have $\text{Aut}(N) = C_{p(p-1)}$. Let $\phi : C_q \rightarrow C_{p(p-1)}$ be the homomorphism defining the semi direct product. In the case that $p$ is odd, we can apply \ref{Classification of semi-direct products between Cyclic p-group and q-groups} to get $C_{p^2 q}$ and if $q \mid p - 1$ then $C_{p^2} \rtimes C_q$. If $p = 2$, then $\phi : C_q \rightarrow C_2$. So only the trivial homomorphism is possible as $q \nmid 2$, giving us $G \cong C_{4q}$.

If $N \cong C_p \times C_p$, then we have $\text{Aut}(N) = GL_2(p)$. TODO

\paragraph{Case 2:} There exists a normal Sylow $q$-group

Let $N$ be the normal Sylow $p$-group. Then $|N| = q$ and by Schur-zassenhaus theorem there exists a complementary subgroup $U$ of $G$ with $|U| = p^2$ such that $G \cong N \rtimes U$. We have $N \cong C_q$ and either $U \cong C_{p^2}$ or $U \cong C_p \times C_p$ (by classification of groups of order $p^2$, TODO - link result)

If $N \cong C_{p^2}$, then again by \ref{Classification of semi-direct products between Cyclic p-group and q-groups} we get these possible groups:
\begin{itemize}
    \item $C_{p^2q}$ always
    \item $C_q \rtimes_1 C_{p^2}$ if $v_p(q - 1) \geq 1$ or $p 
    \mid q - 1$ (Image of action is of order $p$)
    \item $C_q \rtimes_2 C_{p^2}$ if $v_p(q - 1) \geq 2$ or $p^2 \mid q - 1$ (Image of action is of order $p^2$)
\end{itemize}

If $N \cong C_p \times C_p$, then by first isomorphism theorem, $\text{Im}(\phi) \cong (C_p \times C_p) / \ker(\phi)$ so order of image is either $1, p$ or $p^2$. If it is $1$, then we have the trivial homomorphism giving $G \cong C_p \times C_p \times C_{q}$. If it is $p^2$, then we must have $\ker(\phi) = 1$ but $\text{Im}(\phi) \cong C_p \times C_p$ would then be a non-cyclic subgroup of the cyclic group $C_{q-1}$, impossible. If it is $p$, then we get a unique group because TODO

\subsubsection{Reusability and Ease of use of Lemmas}

One of our results regarding cyclic groups was stated in terms of \mintinline{lean4}{CyclicGroup}s. This was fine until we faced a compatibility issue due to the following theorem:

\begin{leancode}
    theorem classify_Cqn_rtimes_Cpm
    {p q : ℕ} [hp : Fact p.Prime] [hq : Fact q.Prime] ...
    (f : CyclicGroup (p ^ m) →* MulAut (CyclicGroup (q ^ n))) :
    ∃! r : ... , Nonempty (SemidirectProduct (CyclicGroup (q ^ n)) (CyclicGroup (p ^ m)) f ≃*
               SemidirectProduct (CyclicGroup (q ^ n)) (CyclicGroup (p ^ m))
                 (canonicalAction ...)))
\end{leancode}

The problem is, \mintinline{lean4}{CyclicGroup q ^ 1} is not definitionally equivalent to \mintinline{lean4}{CyclicGroup q}. The only way to bridge them is an isomorphism which happened to overcomplicate the classification proof. The solution was to create a wrapper that uses \mintinline{lean4}{IsCyclic} and \mintinline{lean4}{Nat.card} to express the same thing but improves reusability a lot as it decouples the code with our own cyclic group implementation. Moreover, now the theorem can be applied to \mintinline{lean4}{CyclicGroup q ^ 1} easily as the $q^1 = q$ calculation happens in \mintinline{lean4}{Nat.card} which is not a problem.

\begin{leancode}
    theorem classify_sdp
    {N K : Type*} [Group N] [Group K] [IsCyclic N] [IsCyclic K]
    {p q : ℕ} [hp : Fact p.Prime] [hq : Fact q.Prime] ...
    (hN : Nat.card N = q ^ n) (hK : Nat.card K = p ^ m)
    (φ : K →* MulAut N) :
    ∃! r : ...,
      Nonempty (SemidirectProduct N K φ ≃*
               SemidirectProduct (CyclicGroup (q ^ n)) (CyclicGroup (p ^ m))
                 (canonicalAction ...)))
\end{leancode}

The final isomorphism still uses our cyclic group and in theory, we could remove that as well but we didn't think it was worth it because you would need to supply the cyclic group and prove its cyclic each time.

\subsection{Order $16$}
\label{Order $16$}

For classifying groups of order $16$, we incorporated the mathematical tooling set out by Marcel Wild \cite{wild_groups_2005}, which is specifically designed to rely solely on elementary group theory facts, and operates by considering the possible cyclic extensions to obtain a group of order 16.
Although such a constructive method may be more verbose, we found that LLM incorporated tools are well suited for generating the required definitions, and can lift most of the burden in the task.

\subsubsection{Overview of Wild's proof}

Here we give an overview of the part of Wild's work which is relevant to our goal, all of which are formalised as a part of our project. Their proofs are temporarily dismissed until we discuss about our implementation in Lean later.

An important machinery Wild defined is the \textit{extension type}.

\begin{definition}[Wild]
    A quadruple $(N,\ n,\ \tau,\ v)$ is an \textit{extension type} if for group $N$, $v \in N$, $n \in \mathbb{N}^+$,  $\tau \in \text{Aut}(N)$, $\tau (v) = v$ and $\tau^n = t_v$ (conjugate by $v$).
\end{definition}

\begin{definition}
    Although not explicitly defined in the text, a group $G$ is said to realise an extension type $(N,\ n,\ \tau,\ v)$ if:
    \begin{itemize}
        \item $N \lhd G$;
        \item There is an inducing element, $a \in G$, such that
        \item $\tau$ is exactly conjugating by $a$ for elements in $N$;
        \item $a^n = v$ (or, equivalently, the coset $Na$ has order $n$ in the quotient group $G/N$);
        \item The quotient group $G/N$ is isomorphic to a cyclic group of order $n$.
    \end{itemize}
\end{definition}

Intuitively, a group realising an extension type means that the group can be constructed from $N$ by layering a cyclic group of order $n$, the two layers interact as directed by $\tau$ and coincides at $v$.

In our project, we refer to $\tau$ by \textit{act} and $v$ by \textit{glue}.

\begin{definition}[Wild]
    Two extension types, $(N,\ n,\ \tau,\ v)$ and $(N',\ n,\ \sigma,\ w)$ are said to be equivalent if there exists isomorphism $\phi : N \isom N'$ such that $\sigma = \phi \circ \tau \circ \phi^{-1}$ and $w = \phi(v)$.
\end{definition}

This is an equivalence relation, and its importance is signified by the following lemma.

\begin{lemma}[Wild]
    Groups $G$ and $G'$ that realise equivalent extension types must be isomorphic.
\end{lemma}

Therefore, by classifying the possible classes of equivalent extension types we can effectively classify the isomorphism classes.

In order to employ this tactic for groups of order 16, further lemmas are established.

\begin{lemma}[Wild]
    Take a group $N$.
    For some automorphic elements $v$ and $w$ in $N$ (i.e. there exists an automorphism which maps $v$ to $w$), and a fixed union of conjugacy classes from $\text{Aut}(N)$, $S$.
    Suppose for $\tau$ ranging over $S$ there is a group $G_\tau$ (and respectively $F_\tau$) which realises the extension type $(N,\ n,\ \tau,\ v)$ (and respectively $(N,\ n,\ \tau,\ w)$) for some $n \in \mathbb{N}$.
    Then the families $\{G_\tau | \tau \in S\}$ and $\{F_\tau | \tau \in S\}$ comprise the same isomorphism classes.
\end{lemma}

\begin{lemma}[Wild]
    Take a group $N$.
    For some characteristic element $v$ in $N$ (i.e. for any automorphism of $N$, $v$ is a fixed point).
    For any two groups $G$ and $G'$ such that they realise some extension types $(N,\ n,\ \tau,\ v)$ and $(N,\ n,\ \sigma,\ v)$ (respectively), if $\sigma$ and $\tau$ belongs to the same conjugacy class in $\text{Aut}(N)$, then $G$ and $G'$ are isomorphic.
\end{lemma}

The above lemmas significantly reduce the choices of equivalent extension type classes for a given base group $N$. This enables the possibility to perform an exhaustive classification across all equivalence classes. However, to apply to groups of order 16, a further lemma is needed to restrict the choice of possible $N$'s.

\begin{lemma}[Wild]
    If $G$ is a group of order $16$ that is not isomorphic to $C_2 \times C_2 \times C_2 \times C_2$, then either $C_8 \lhd G$ or $C_4 \times C_2 \lhd G$.
\end{lemma}

Using the above lemmas, it is now possible to do a by case analysis to realise the possible equivalence classes, the main result is stated below.

\begin{lemma}[Wild]
    Lemma 3
\end{lemma}

\subsection{Order $24$}
\label{Order $24$}

\subsection{Order $30$}
\label{Order $30$}

We follow the general structure of the proof given in the course handout by Deopurkar \cite{deopurkar_modern_2013}. We have generalised some lemmas used in the proof to increase its usefulness.

\begin{theorem}\label{Groups of order $pqr$ have either a normal Sylow-$r$ group or normal Sylow-$q$ group}
    A group $G$ of order $pqr$ for distinct primes $p < q < r$ must have either a normal Sylow-$r$ group or normal Sylow-$q$ group.
\end{theorem}
\begin{proof}
    By Sylow's theorems, we get $n_r \equiv 1 \pmod r$ and $n_r \mid pq$. If $n_r = 1$  we are done. $n_r = p$ or $q$ cannot hold because $p, q < r$ and $n_r \equiv 1 \pmod r$. Thus $n_r = pq$, and there are $pq(r-1)$ elements of order $r$. We have $n_q \mid pr, n_q \equiv 1 \pmod q$. Thus if $n_q > 1$, then $n_q \geq q + 1 > p + 1$. Hence, there are atleast $(p+1)(q-1) = pq + q - p - 1 \geq pq$ elements of order $q$, a contradiction as we already have $pqr$ elements but the identity has order $1$. Hence, $n_q = 1$ and we are done.
\end{proof}

To show the next theorem, we use the following well known results:
\begin{lemma}\label{Product of subgroups is subgroup if one is normal}
    Let $N$ and $H$ be subgroups of $G$ with $N \lhd G$. Then $NH$ is a subgroup of $G$.
\end{lemma}
\begin{lemma}\label{Order of product of subgroups}
    For subgroups $H, K$ of $G$, we must have $|HK| = \frac{|H||K|}{|H \cap K|}$
\end{lemma}
\begin{lemma}\label{Subgroup of index smallest prime is normal}
    Let $G$ be a finite group and $p$ be the smallest prime dividing $|G|$. Then any subgroup of $G$ of index $p$ is normal.
\end{lemma}

\begin{theorem}\label{Group of order $pqr$ has a normal subgroup of order $qr$}
    A group $G$ of order $pqr$ for distinct primes $p < q < r$ has a normal subgroup of order $qr$
\end{theorem}
\begin{proof}
    By \ref{Groups of order $pqr$ have either a normal Sylow-$r$ group or normal Sylow-$q$ group}, we either have a normal Sylow-$r$ group or normal Sylow-$q$ group. If we have a normal Sylow-$r$ group $N$, then $NH$ is a subgroup (by \ref{Product of subgroups is subgroup if one is normal}) of $G$ of order $qr$ (by \ref{Order of product of subgroups} and the fact that groups of prime order are trivially intersecting) where $H$ is any Sylow-$q$ group. Then $NH \lhd G$ follows from \ref{Subgroup of index smallest prime is normal}. An analogous argument holds for the case when there is a normal Sylow-$q$ group.
\end{proof}

We are now ready to present the main proof (adapted from the course handout by Deopurkar \cite{deopurkar_modern_2013}).

Let $G$ be a group of order $30$. Let $N \lhd G$ be of order $15$, which exists by \ref{Group of order $pqr$ has a normal subgroup of order $qr$}.
By Schur-zassenhaus theorem there exists a complementary subgroup $U$ of $G$ with $|U| = 2$ such that $G \cong N \rtimes U$. We have $U \cong C_2$ and by the classification of groups of order $pq$ we have $N \cong C_{15}$. 

The underlying homomorphism of the semi-direct product is $\phi : C_2 \rightarrow \text{Aut}(C_{15})$. The image of the homomorphism must have order $1$ or $2$. If the image order is $1$, then $G \cong C_{15} \times C_2 \cong C_{30}$. 

Now suppose image has order $2$. By \ref{$K$ cyclic $p$-group then equal images define isomorphic semi-direct products}, the group $G$ is isomorphic to the semi-direct product defined by $\phi$ whose image is generated by a particular element of order $2$ of $\text{Aut}(C_{15})$. Now $\text{Aut}(C_{15}) \cong (\mathbb{Z}/15\mathbb{Z})^\times$ has exactly $3$ elements of order $2$ (namely $4, 11, 14$), which gives us $3$ non-trivial semi-direct products, giving us a total of $4$ groups. A nice presentation of the groups is: $C_{30}, D_{15}, C_5 \times D_3, C_3 \times D_5$.

\section{Uniqueness}

\section{AI}
We used AI to write Lean code.
We used a few different AI: Claude, Claude code, Gemini, Aristotle, Numina Fuse.

\section{Why we cannot do 32}
The invariant strategy we are using for uniqueness breaks.

% Kevin: the great thing about Lean is that you only have to read the statements
% so you can get an AI to do the proofs and not read these

%%%%%%%%%%%%%%%%%%%%%%%%%%%%%%%%%%%%%
%%%%%%%%%%%%%%%%%%%%%%%%%%%%%%%%%%%%%
% Acknowledgements
%%%%%%%%%%%%%%%%%%%%%%%%%%%%%%%%%%%%%
\section*{Acknowledgments}
\addcontentsline{toc}{section}{Acknowledgement}
Katerina Hristova

Xena Club

Bhavik Mehta

Edison Xie

London Learning Lean

Neil Bhat
%%%%%%%%%%%%%%%%%%%%%%%%%%%%%%%%%%%%%

%%%%%%%%%%%%%%%%%%%%%%%%%%%%%%%%%%%%%
% Appendices
%%%%%%%%%%%%%%%%%%%%%%%%%%%%%%%%%%%%%
\appendix

\section{First appendix}
\label{sec:appendix1}

\section{Second appendix}
\label{sec:appendix2}

%%%%%%%%%%%%%%%%%%%%%%%%%%%%%%%%%%%%%

%%%%%%%%%%%%%%%%%%%%%%%%%%%%%%%%%%%%%
\bibliography{bibliography.bib}
%%%%%%%%%%%%%%%%%%%%%%%%%%%%%%%%%%%%%

\end{document}

%%FILE numina/blueprints/redoing-pq-classification-without-order-pq-library/redoing-pq-classification-without-order-pq-library.tex
% Blueprint: Redoing pq classification without order pq library
% Created by Numina

\begin{document}

% Your blueprint will go here.

\end{document}


%%FILE numina/blueprints/redoing-pq-classification-without-order-pq-library/source/redoing-pq-classification-without-order-pq-library-source.tex
Use the computable canonicalAction in CycPGroupClassification.lean and classify_Cqn_rtimes_Cpm_exists or classify_Cqn_rtimes_Cpm not sure which one. Take inspiration from 4q and 2p^2 classifications. pq should be quite simple though. The proof is:

(Show this as a lemma in SylowUtils.lean, shouldn't be more than like 50 lines probably)
\begin{lemma}
    Let $G$ be a group of order $pq$ for primes $p < q$. Then $G$ has a normal Sylow-$q$ subgroup.
\end{lemma}
\begin{proof}
    By Sylow's theorems, we get $n_q \equiv 1 \pmod q$ and $n_q \mid p$. If $n_q = p$, then as $p < q$, we we get a contradiction to $n_q \equiv 1 \pmod q$. Hence $n_q = 1$ and so $G$ has a normal Sylow-$q$ subgroup.
\end{proof}

Then: A corollary of the lemma is that $G \cong C_q \rtimes C_p$. We have already proven a much more general lemma \ref{Classification of semidirect products between Cyclic p-group and q-groups} which shows that $G$ is isomorphic to either $C_{pq}$ or the non-trivial semi-direct product $C_q \rtimes C_p$ where the image of the action $f : C_p \rightarrow C_{q - 1}$ has order $p$ when we have $p \mid q - 1$.
(The \ref{Classification of semidirect products between Cyclic p-group and q-groups} is basically canonicalAction and classify_Cqn_rtimes_Cpm_exists etc.)

Please read all files under Lemmas/ directory first so you don't reprove anything. Currently, the project uses this OrderPQ library to do the pq classification and the task I am giving you is to basically do the classification ourselves, just like 4q and 2p^2 are completely proven in-house. This entails a similar classification theorem to 4q and 2p^2. You should write p < q and then condition for second one would be whatever p divides q - 1 (just like we have conditions for some 4q groups). Then you need to make sure the canonicalAction groups are still being used in the small groups library properly and in the Classification.lean.

Good luck! Should be a quick ride.

%%FILE numina/blueprints/refactoring-retrieve-to-make-it-computable/refactoring-retrieve-to-make-it-computable.tex
% Blueprint: Refactoring retrieve to make it computable
% Created by Numina

\begin{document}

\section{Goal of this blueprint}

We are classifying finite groups $G$ of order $p^2 q$ with $p \neq q$ primes,
and in particular the $4q$ case ($p = 2$, $q \geq 3$ odd prime). The Lean code
has accumulated ad-hoc bridging logic in the $4q$ case that re-proves things
which are already established by the general classification machinery
(\texttt{classify\_Cqn\_rtimes\_Cpm}) living in
\texttt{M2rGroup7.P2qClassification.CycPGroupClassification}. The goal of this
refactor is to make \texttt{classify\_Cqn\_rtimes\_Cpm} the single source of
truth wherever the math proof says ``use Lemma~3''.

This blueprint records the informal mathematical proof verbatim, then states
how each step maps to existing Lean machinery. The refactor target is
\texttt{classification\_4q} in
\texttt{M2rGroup7/P2qClassification/FourQClassification.lean}.

\section{Lemmas}

\begin{lemma}[Normal Sylow subgroup]
    \label{lem:p2q-normal-sylow}
    For groups of order $p^2 q$, there must exist either a normal Sylow
    $p$-group or a normal Sylow $q$-group.
\end{lemma}

\begin{lemma}[Same image in $\operatorname{Aut}(H)$ gives isomorphic SDPs]
    \label{lem:same-image-sdp}
    If $K$ is a cyclic $p$-group, two homomorphisms $f, g : K \to
    \operatorname{Aut}(H)$ define isomorphic semi-direct products $H \rtimes K$
    if they have the same image.
\end{lemma}

\begin{lemma}[Classification of $C_{q^n} \rtimes C_{p^m}$]
    \label{lem:classify-Cqn-Cpm}
    Semi-direct products $C_{q^n} \rtimes C_{p^m}$ ($q$ odd prime, $p \neq q$)
    are classified up to isomorphism by $r \in \{0, \ldots, \min(m, d)\}$
    where $d = v_p(q - 1)$, giving $\min(m, d) + 1$ classes. Every action $f$
    belongs to exactly one class.
\end{lemma}

\begin{lemma}[Conjugate images give isomorphic SDPs]
    \label{lem:conjugate-image-sdp}
    If $|K| = p$, then two homomorphisms $f, g : K \to \operatorname{Aut}(H)$
    define isomorphic semi-direct products $H \rtimes K$ if their images are
    conjugate subgroups of $\operatorname{Aut}(H)$.
\end{lemma}

\section{Case 1: a normal Sylow $p$-subgroup exists}

$|N| = p^2$ and by Schur--Zassenhaus there exists a complementary subgroup $U$
of $G$ with $|U| = q$ such that $G \cong N \rtimes U$. We have $U \cong C_q$.

\subsection{Case 1.1: $N \cong C_{p^2}$}

$\operatorname{Aut}(N) \cong C_{p(p-1)}$. Let
$\varphi : C_q \to C_{p(p-1)}$. Done by Lemma~\ref{lem:classify-Cqn-Cpm}
when $p$ is odd. When $p = 2$, we get $\varphi : C_q \to C_2$ so only the
trivial homomorphism is possible, as $q \nmid 2$ for odd $q$.

\begin{itemize}
    \item (For $p = 2$, only trivial homomorphism.)
    \item (For $q = 2$, we get $2$ homomorphisms.)
\end{itemize}

Outcomes:
\begin{itemize}
    \item $C_{p^2 q}$ always.
    \item $C_{p^2} \rtimes C_q$ if $q \mid p - 1$.
\end{itemize}

\subsection{Case 1.2: $N \cong C_p \times C_p$}

$\operatorname{Aut}(N) \cong \operatorname{GL}_2(p)$. Let
$\varphi : C_q \to \operatorname{GL}_2(p)$.

\paragraph{Special case $p = 2$.}
$\operatorname{GL}_2(2) \cong S_3$. Either $\varphi$ is a trivial action, or
its image is a subgroup of order $q$ in $S_3$, hence $q = 3$. There is only
one subgroup of order $3$, so this gives one group.

\paragraph{Special case $q = 2$.}
$\operatorname{Im}(\varphi)$ is trivial or $\langle A \rangle$ where $A \in
\operatorname{GL}_2(p)$ is of order $2$. By Lemma~\ref{lem:conjugate-image-sdp}
we only need to look at similarity classes and pick a representative from
each class. Then $A^2 = 1$ implies the minimal polynomial of $A$ divides
$x^2 - 1$, and since $p$ is odd, it has distinct factors and so $A$ is
diagonalisable. Each such matrix is similar to a diagonal matrix
$\operatorname{diag}(\alpha, \beta)$ with $\alpha^2 = \beta^2 = 1$, hence
$\alpha = \pm 1$, $\beta = \pm 1$ (since $t^2 - 1$ has at most $2$ roots over a
field). The case $\alpha = \beta = 1$ has order $1$, so we discard it. There
are two similarity classes: $\alpha = 1, \beta = -1$ (equivalently
$\beta = 1, \alpha = -1$), and $\alpha = \beta = -1$.

\emph{Side note:} Lemma~\ref{lem:conjugate-image-sdp} is not strictly needed
here because the general condition $\varphi_2(u) = \alpha \circ
\varphi_1(\beta(u)) \circ \alpha^{-1}$ becomes equivalent to similarity:
$\operatorname{Aut}(C_q) = C_1$ when $q = 2$ so $\beta = \operatorname{id}$,
and $\varphi_1, \varphi_2$ are determined by where they send the generator
since $U$ is cyclic.

\paragraph{General special case.} A degree-$3$ map $C_3 \to
\operatorname{Aut}(G)$, when applicable.

\section{Case 2: a normal Sylow $q$-subgroup exists}

$|N| = q$ and by Schur--Zassenhaus there exists a complementary subgroup $U$
of $G$ with $|U| = p^2$ such that $G \cong N \rtimes U$. We have $N \cong
C_q$ and $\operatorname{Aut}(N) \cong C_{q-1}$.

\subsection{Case 2.1: $U \cong C_{p^2}$}

Let $\varphi : C_{p^2} \to C_{q-1}$. Done by
Lemma~\ref{lem:classify-Cqn-Cpm} when $q$ is odd. Only the trivial
homomorphism is possible when $q = 2$.

\begin{itemize}
    \item (When $p = 2$, have to see if $q$ is $1$ mod $4$.)
    \item (When $q = 2$, we get only the trivial homomorphism.)
\end{itemize}

Outcomes:
\begin{itemize}
    \item $C_{p^2 q}$ always.
    \item $C_q \rtimes_1 C_{p^2}$ if $v_p(q - 1) \geq 1$.
    \item $C_q \rtimes_2 C_{p^2}$ if $v_p(q - 1) \geq 2$.
\end{itemize}

\subsection{Case 2.2: $U \cong C_p \times C_p$}

Let $\varphi : C_p \times C_p \to C_{q-1}$.

$\operatorname{Im}(\varphi) \cong (C_p \times C_p) / \ker(\varphi)$, so the
order of the image is either $1$, $p$, or $p^2$.
\begin{itemize}
    \item If it is $1$: $\varphi$ is trivial, giving
        $G \cong C_p \times C_p \times C_q$.
    \item If it is $p^2$: $\ker(\varphi) = 1$, so $\operatorname{Im}(\varphi)$
        is a non-cyclic subgroup of the cyclic group $C_{q-1}$, impossible.
    \item If it is $p$: ad-hoc argument matching $\varphi$ against the
        canonical action of image order $p$ (sub-case of
        Lemma~\ref{lem:same-image-sdp} applied with cyclic codomain
        $\operatorname{Aut}(C_q)$, where same image cardinality implies same
        image because cyclic groups have a unique subgroup of any given order).
\end{itemize}

\begin{itemize}
    \item (When $p = 2$, one trivial and one non-trivial.)
    \item (When $q = 2$, we get only the trivial homomorphism.)
\end{itemize}

Outcomes:
\begin{itemize}
    \item $C_p \times C_p \times C_q$.
    \item $C_q \rtimes (C_p \times C_p)$ if $p \mid q - 1$.
\end{itemize}

\section{Mapping to Lean machinery}

The mapping from the math sketch to existing Lean lemmas:

\begin{itemize}
    \item \textbf{Lemma~\ref{lem:p2q-normal-sylow}} $\leftrightarrow$
        \texttt{p2q\_group\_has\_normal\_sylow\_subgroup} in
        \texttt{M2rGroup7.Lemmas.SylowUtils}.
    \item \textbf{Lemma~\ref{lem:same-image-sdp}} $\leftrightarrow$
        \texttt{semidirectProduct\_iso\_if\_range\_eq} (and its cardinality
        variant \texttt{semidirectProduct\_iso\_if\_range\_card\_eq}, which
        upgrades equal-cardinality to equal-image using uniqueness of cyclic
        subgroups in a cyclic ambient group) in
        \texttt{M2rGroup7.Lemmas.ClassificationUtils}.
    \item \textbf{Lemma~\ref{lem:classify-Cqn-Cpm}} $\leftrightarrow$
        \texttt{classify\_Cqn\_rtimes\_Cpm} (and the explicit-$r$ variant
        \texttt{classify\_Cqn\_rtimes\_Cpm\_exists}) in
        \texttt{M2rGroup7.P2qClassification.CycPGroupClassification}.
    \item \textbf{Lemma~\ref{lem:conjugate-image-sdp}} is not currently
        formalised at this level of generality; the only place it is used in
        the math sketch above is Case~1.2 $q = 2$, which the
        \texttt{4q} classification does not enter (we have $q \geq 3$). The
        side note explains why the lemma is unnecessary in that case anyway.
\end{itemize}

For $4q$ specifically (i.e.\ $p = 2$, $q \geq 3$), Case~1.1 needs the special
$p = 2$ argument (no use of Lemma~\ref{lem:classify-Cqn-Cpm}); Case~1.2 enters
the special $p = 2$ sub-case, which gives the $C_3 \to \operatorname{Aut}(C_2
\times C_2)$ branch (only when $q = 3$); Case~2.1 directly invokes
Lemma~\ref{lem:classify-Cqn-Cpm} with parameters $p = 2$, $m = 2$, $q$ odd,
$n = 1$; Case~2.2 follows the ad-hoc ``image trivial or of order $p = 2$''
argument and uses \texttt{semidirectProduct\_CpCp\_iso}.

\section{Refactor target}

\paragraph{Current state.} In
\texttt{M2rGroup7/P2qClassification/FourQClassification.lean}, the Case~2.1
($C_4$) branch of \texttt{classification\_4q} currently dispatches through an
intermediate helper \texttt{classify\_sdp\_C4\_on\_Cq} that re-proves the
content of Lemma~\ref{lem:classify-Cqn-Cpm} via range-card matching. This
introduces auxiliary lemmas
(\texttt{canonicalC4OnCqAction\_range\_card},
\texttt{canonicalC4OnCqAction\_r2\_range\_card}) which exist only to feed
\texttt{semidirectProduct\_iso\_if\_range\_card\_eq}.

\paragraph{Target.} The Case~2.1 branch should invoke
\texttt{classify\_Cqn\_rtimes\_Cpm} directly on $\varphi$ (transported to
$C_{2^2} \to \operatorname{Aut}(C_{q^1})$), then transport the resulting
isomorphism back to $C_4 \to \operatorname{Aut}(C_q)$ form. The transport
between $C_{2^2}$ and $C_4$, and between $C_{q^1}$ and $C_q$, is genuine
Lean-only bookkeeping (definitional equality on the underlying types is
prevented by the indexing) and is handled by a single semi-direct-product
transport lemma \texttt{SemidirectProduct.transportCpCqIso} together with
\texttt{transportCpCqHom} (already present). The matching to
\texttt{canonicalC4OnCqAction}, \texttt{canonicalC4OnCqAction\_r2} is then by
definitional equality since these wrappers are themselves defined as
\texttt{transportCpCqHom} of the appropriate \texttt{canonicalAction}.

\paragraph{What stays.} The Case~2.2 ($C_2 \times C_2$) branch is genuinely
ad-hoc per the math sketch and remains supported by
\texttt{classify\_sdp\_C2C2\_on\_Cq} plus the single range-card lemma
\texttt{canonicalC2C2OnCqAction\_range\_card}. The Case~1.2 ($q = 3$) sub-case
likewise remains ad-hoc and uses \texttt{canonicalC3OnC2C2Action\_range\_card}
via \texttt{semidirectProduct\_C3\_on\_C2C2\_iso}.

\paragraph{What goes.} \texttt{classify\_sdp\_C4\_on\_Cq},
\texttt{canonicalC4OnCqAction\_range\_card}, and
\texttt{canonicalC4OnCqAction\_r2\_range\_card} are deleted from
\texttt{FourQClassification.lean}; they re-prove
Lemma~\ref{lem:classify-Cqn-Cpm}'s content rather than invoking it.

\end{document}
% Blueprint: Refactoring retrieve to make it computable
% Created by Numina

\begin{document}

% Your blueprint will go here.

\end{document}


%%FILE numina/blueprints/refactoring-retrieve-to-make-it-computable/source/refactoring-retrieve-to-make-it-computable-source.tex
You have to figure it out. Some things Claude wrote but do your own research:

"Now I have the full picture. Let me correct and sharpen my earlier analysis.

  The actual noncomputability chain:

  zmodAddCyclicAddEquiv  ← uses h.exists_generator.choose  [root cause]
    ↑
  zmodCyclicMulEquiv     [noncomputable]
    ↑
  mulEquivOfCyclicCardEq [noncomputable]  ← the Mathlib lemma you're pointing to
    ↑
  sdpCanonicalAction     (uses it for eN, eK)  [noncomputable]
    ↑
  canonicalCpOnCqAction  [noncomputable]

  AND separately in canonicalAutElement_exists:
  IsCyclic.exists_generator (α := (ZMod q^n)ˣ).choose  [noncomputable]
  IsCyclic.mulAutMulEquiv                                [noncomputable, goes through zmodCyclicMulEquiv]

  ---
  Can mulEquivOfCyclicCardEq be made computable?
  
  The Mathlib version is noncomputable because it works for an arbitrary abstract cyclic group G — it must pick a generator nonconstructively. But for our concrete CyclicGroup n = Multiplicative (ZMod n),
  the generator Multiplicative.ofAdd 1 is explicitly known. Mathlib even provides the explicit-generator version:

  zmodMulEquivOfGenerator  -- takes an explicit generator, still marked noncomputable

  The issue is that even with an explicit generator, zmodAddEquivOfGenerator goes through QuotientAddGroup.liftEquiv which involves quotient type machinery that Lean marks noncomputable regardless. So the
  noncomputability can't be lifted just by providing a generator.

  However, for the specific case CyclicGroup n ≃* CyclicGroup n (same type), it's just MulEquiv.refl — trivially computable, no iso machinery needed. This covers the case in sdpCanonicalAction with n=m=1,
  where eN : CyclicGroup q ≃* CyclicGroup (q^1) — once you handle the q = q^1 nat equality, it collapses to refl.

  ---
  Can canonicalAutElement be made computable after that?
  
  Even if you fix mulEquivOfCyclicCardEq, canonicalAutElement has its own noncomputable .choose finding a generator of (ZMod q^n)ˣ. But since (ZMod q)ˣ is Fintype and DecidableEq, finding a generator is 
  computationally decidable — you can search the finite set for an element with orderOf g = q-1. For instance, Finset.filter (fun g => orderOf g = q - 1) Finset.univ is nonempty, and you can take .choose on
  that — but at this point it's a decidable choose, which Lean can compute.

  More concretely for our r=1 case: you just need a unit u ∈ (ZMod q)ˣ with u^p = 1 and u ≠ 1, and the corresponding "multiply by u" automorphism. Since (ZMod q)ˣ is finite and the predicate is decidable,
  Finset.choose gives a computable element. And the automorphism Multiplicative.ofAdd a ↦ Multiplicative.ofAdd (u * a) is concretely computable.

  So yes, in principle both are fixable, but it requires concrete refactoring:
  1. Replace mulEquivOfCyclicCardEq in sdpCanonicalAction with a computable identity iso for the CyclicGroup case
  2. Replace IsCyclic.exists_generator.choose in canonicalAutElement with a decidable Finset.choose searching (ZMod q)ˣ for a primitive root
  3. Replace IsCyclic.mulAutMulEquiv with a computable concrete iso MulAut (CyclicGroup q) ≃* (ZMod q)ˣ via f ↦ f (Multiplicative.ofAdd 1)
  
  The practical difficulty is that the CycPGroupClassification.lean proofs are written abstractly for any cyclic group — making them computable would mean either specializing everything to CyclicGroup n
  concretely, or finding computable analogues of each abstract step."
